\chapter{Complex Hessian equations}
\label{chap:hessian-equations}

Let $X$ be a connected compact complex manifold of dimension $n$,
let $\omega$ be a smooth positive Hermitian form, and fix $2\leq m<n$.
The complex Hessian operator
$(\omega+\ddc u)^m\wedge\omega^{n-m}$ replaces the determinant by an
intermediate symmetric polynomial. Its admissible cone allows negative
eigenvalues, so the product calculus for plurisubharmonic potentials
from \cref{chap:framework} no longer applies directly.

For a continuous admissible weak solution with smooth positive density
depending on the potential, the difficulty is to obtain derivative
estimates before smoothness is known. Proper exponential equations
associate smooth solutions to its smooth admissible approximations.
The complex-Hessian estimate is independent of derivatives of the
prescribed approximation, so it survives passage to the limit.
Uniform convergence defines the continuous mixed products, while
obstacle envelopes give the domination principle that identifies the
limit with the original solution. A concave uniformly elliptic extension
then supplies classical interior regularity. The determinant case is treated in
\cref{thm:bounded-hermitian-smoothness}; the argument here concerns the
intermediate Hessian degrees.
Uniform stability also constructs solutions for strictly increasing
exponents, by contraction between constant barriers.

Bounded admissible potentials require decreasing approximation. The
recursion first gives currents with measure coefficients below degree
$m$; local Hessian capacities control multiplication by their
discontinuous potentials. Quasicontinuity and weighted convergence
make the products symmetric and allow all approximation indices to
tend to their limits together. At degree $m$, the positive product
after wedging with $\omega^{n-m}$ is a Radon measure.

These weighted limits pass locality and the maximum inequality to
bounded potentials. Truncation above the infimum then proves that
the Hessian measure of each potential is nonzero on an arbitrary
Hermitian background. In global comparison, the maximum inequality
leaves a signed total-mass difference. The complex Hessians of the
background powers determine when this difference vanishes for every
bounded admissible list. On small sublevel sets, a modified comparison
estimate absorbs the torsion error on arbitrary Hermitian backgrounds.
It yields strict domination, bounded stability for proper equations and
smoothness for their solutions with smooth data between constant barriers.

\section{The admissible cone}
\label{sec:hessian-cone}
\index{Hessian equation!admissible cone}\index{Garding cone@G\aa rding cone}

For $\lambda\in\mathbb R^n$, let $S_j(\lambda)$ be the elementary
symmetric polynomial of degree $j$, put $\sigma_j=S_j/\binom nj$ and
$\sigma_0=1$, and set
\[
 \Gamma_m=\{\lambda:S_j(\lambda)>0,\ 1\leq j\leq m\}.
\]
A real $(1,1)$-form belongs to $\Gamma_m(\omega)$ if its eigenvalues
relative to $\omega$ belong to $\Gamma_m$. Write $C_{n,m}$ for the
corresponding closed cone of Hermitian matrices. The algebraic assertions
below also apply when $m=n$.

We use the classical G\aa rding facts: the cone $C_{n,m}$ is convex,
contains the nonnegative matrices, and is preserved by adding them;
adding a positive matrix puts a matrix of $C_{n,m}$ in its interior.
The normalized root $F=\sigma_m^{1/m}$ is concave and increasing on
$\Gamma_m$. Polarized products of $k$ forms in $\Gamma_k(\omega)$ with
$\omega^{n-k}$ are positive, and on $\Gamma_k$, hence by continuity on $\overline{\Gamma_k}$, Maclaurin's inequalities
give $\sigma_j^{1/j}\geq\sigma_k^{1/k}$ for $1\leq j\leq k$. If $\alpha$
has eigenvalues $\lambda$ relative to $\omega$, then
$\alpha^j\wedge\omega^{n-j}=\sigma_j(\lambda)\,\omega^n$; this is the
normalization of all Hessian measures in the chapter.

A continuous real function $v$ on $X$ is \emph{admissible} if
\begin{equation}\label{eq:hessian-continuous-admissibility}
 (\omega+\ddc v)\wedge\gamma_1\wedge\cdots\wedge\gamma_{m-1}
                                  \wedge\omega^{n-m}\geq0
\end{equation}
in distributions for every smooth collection
$\gamma_1,\ldots,\gamma_{m-1}\in\Gamma_m(\omega)$.
For smooth $v$ this is equivalent, by G\aa rding duality, to
$\omega+\ddc v\in\overline{\Gamma_m(\omega)}$ at every point; a smooth
$v$ is \emph{strictly admissible} if $\omega+\ddc v\in\Gamma_m(\omega)$
everywhere. Maxima and convex combinations preserve continuous
admissibility. To check the maximum property, fix
$\gamma_1,\ldots,\gamma_{m-1}$ in
\eqref{eq:hessian-continuous-admissibility}. By polarized positivity,
the left-hand side is $(Lv+g)\,\omega^n$ for a linear elliptic operator
$L$ with smooth coefficients and a smooth function $g$. Subtracting a
local smooth solution of $Lh=-g$ reduces the assertion to the classical
fact that the maximum of two continuous distributional subsolutions of
$L$ is again a subsolution.

For a bounded upper semicontinuous function $v$, admissibility means that
every \emph{upper test} $q$ at a point $x$, that is, every $C^2$ function
$q$ near $x$ with $q\geq v$ near $x$ and $q(x)=v(x)$, satisfies
\begin{equation}\label{eq:hessian-bounded-upper-tests}
 \omega(x)+\ddc q(x)\in\overline{\Gamma_m(\omega(x))}.
\end{equation}
For continuous $v$, this agrees with
\eqref{eq:hessian-continuous-admissibility}: G\aa rding duality reduces
both conditions to subsolution inequalities for smooth strictly elliptic
linear operators, and we use the classical equivalence of viscosity and
distributional subsolutions for such operators. In the bounded class,
pairing \eqref{eq:hessian-bounded-upper-tests} with $\omega^{n-1}$
makes $v$, up to a local smooth function, a viscosity subsolution, and
hence a subharmonic function, for such an operator. As for ordinary
subharmonic functions, $v(x)$ is then the limit of the essential
suprema of $v$ on small coordinate balls centered at $x$. Thus $v$ is
the \emph{canonical representative} of its class in $L^1(X)$, and
this representative is determined by the distribution defined by $v$.
\index{Admissible potential!bounded upper tests}\index{Admissible potential!canonical representative}

The first estimate controls a scalar shift of the cone. It will give
lower bounds for every coefficient of the linearized Hessian operator.

\Needspace{14\baselineskip}
\begin{lemma}\label{hessian:hessian-cone-shift}
For $\lambda\in\Gamma_m$, put
\[
 e_m=\frac{\sigma_m(\lambda)}{m\sigma_{m-1}(\lambda)}.
\]
Then $\lambda-t\mathbf1\in\Gamma_m$ for $0\leq t<e_m$. If $m\geq2$, the conclusion also holds for $t=e_m$. For $m=1$, the endpoint lies on the boundary.
\end{lemma}
\begin{proof}
The polynomial
\[
 P(x)=\sigma_m(\lambda+x\mathbf1)
     =\sum_{j=0}^m\binom mj\sigma_j(\lambda)x^{m-j}
\]
is a positive multiple of the $(n-m)$th derivative of $\prod_i(x+\lambda_i)$. Hence all its roots are real, by Rolle's theorem. Its coefficients are positive, so all its roots are strictly negative. Write $P(x)=\prod_{i=1}^m(x+r_i)$ with $r_i>0$. Logarithmic differentiation at zero gives
\[
 \sum_{i=1}^m r_i^{-1}
       =\frac{m\sigma_{m-1}(\lambda)}{\sigma_m(\lambda)}
       =e_m^{-1}.
\]
Thus $\min_i r_i\geq e_m$, with strict inequality if $m\geq2$. For $x> -\min_i r_i$, every derivative of $P$ is positive: differentiating its product expansion produces sums of products of positive factors. Since
\[
 P^{(m-j)}(x)=\frac{m!}{j!}\sigma_j(\lambda+x\mathbf1),
 \qquad 1\leq j\leq m,
\]
all the required symmetric polynomials are positive at $x=-t$. When $m=1$, the sole root is $-e_1$.
\end{proof}

\Needspace{10\baselineskip}
\begin{lemma}\label{hessian:hessian-mixed-trace}
If $\alpha,\gamma\in\Gamma_m(\omega)$, then
\[
 \frac{\alpha^{m-1}\wedge\gamma\wedge\omega^{n-m}}{\omega^n}
 \frac{\alpha\wedge\omega^{n-1}}{\omega^n}
 \geq\frac1{m!}
 \frac{\alpha^m\wedge\omega^{n-m}}{\omega^n}
 \frac{\gamma\wedge\omega^{n-1}}{\omega^n}.
\]
The assertion extends to $\gamma\in\overline{\Gamma_m(\omega)}$ by continuity.
\end{lemma}
\begin{proof}
Write $\sigma_k=\alpha^k\wedge\omega^{n-k}/\omega^n$, and
$M_k=\alpha^{k-1}\wedge\gamma\wedge\omega^{n-k}/\omega^n$.
Let $1\leq k<m$. For $0\leq t<\sigma_{k+1}/((k+1)\sigma_k)$, \cref{hessian:hessian-cone-shift}, with $k+1$ in place of $m$, puts $\alpha-t\omega$ in $\Gamma_{k+1}(\omega)$. G\aa rding positivity for the $k+1$ forms $\alpha-t\omega,\alpha,\ldots,\alpha,\gamma$ of $\Gamma_{k+1}(\omega)$ gives $(\alpha-t\omega)\wedge\alpha^{k-1}\wedge\gamma\wedge\omega^{n-k-1}\geq0$, that is,
\[
 M_{k+1}\geq tM_k
 \quad\left(0\leq t<\frac{\sigma_{k+1}}{(k+1)\sigma_k}\right).
\]
Let $t$ increase to the displayed endpoint. Iterating, starting with $M_1=\gamma\wedge\omega^{n-1}/\omega^n$, and using $\prod_{k=1}^{m-1}\sigma_{k+1}/((k+1)\sigma_k)=\sigma_m/(m!\sigma_1)$, yields
\[
 M_m\geq\frac{\sigma_m}{m!\sigma_1}M_1.
\]
All denominators are positive. Approximating a form in the closed cone by $\gamma+\delta\omega$, $\delta\downarrow0$, gives the last assertion.
\end{proof}

At a diagonal Hermitian matrix $\Lambda$ in the interior of $C_{n,m}$, order its
eigenvalues as $\lambda_1\geq\cdots\geq\lambda_n$ and write
\[
 f=F(\Lambda),\quad F_i=\partial_{\lambda_i}F,\quad
 S=\operatorname{tr}\Lambda,\quad \mathcal F=\sum_iF_i,
 \quad Q=\sum_iF_i\lambda_i^2.
\]
The next inequalities use these derivatives of the normalized root,
rather than derivatives of the polynomial $S_m$.

\Needspace{10\baselineskip}
\begin{lemma}\label{hessian:weighted-cone-trace}
If $M\in C_{n,m}$, then
\[
 W_M:=\sum_iF_iM_{i\bar i}\geq\frac{f}{m!S}\operatorname{tr}M,
 \qquad F_i\geq\frac{f}{m!S},\qquad
 |M_{1\bar1}|\leq\operatorname{tr}M.
\]
\end{lemma}
\begin{proof}
In the notation of the proof of \cref{hessian:hessian-mixed-trace}, with $\alpha$ and $\gamma$ the forms with matrices $\Lambda$ and $M$, differentiation gives $DF_\Lambda(M)=fM_m/\sigma_m$. The estimate $M_m\geq\sigma_m M_1/(m!\sigma_1)$, with $M_1=\operatorname{tr}M/n$ and $\sigma_1=S/n$, proves the first inequality. The positive semidefinite rank-one coordinate projections belong to $C_{n,m}$; taking $M$ to be each projection proves the second inequality.
Since $m\geq2$, $C_{n,m}$ is contained in $C_{n,2}$, so the eigenvalues $\mu$ of $M$ satisfy
\[
 \operatorname{tr}M\geq0,\qquad
 \sum_i\mu_i^2=(\operatorname{tr}M)^2-2S_2(\mu)
                   \leq(\operatorname{tr}M)^2.
\]
Thus the operator norm of $M$ is at most its trace, which proves the last inequality. The wedge pairing computes the derivative in the diagonal frame of $\Lambda$, regardless of the eigenvectors of $M$.
\end{proof}

\begin{lemma}\label{hessian:quadratic-coercivity}
If $0<f_0\leq f\leq f_1$, then
\[
 Q:=\sum_iF_i\lambda_i^2
       \geq\frac{f_0}{m}S-\frac{n-m}{m}f_1^2,
 \qquad \mathcal F\geq1.
\]
\end{lemma}
\begin{proof}
Write $(\lambda|i)$ for $\lambda$ with its $i$th coordinate omitted. The derivative of the normalized root is $F_i=fS_{m-1}(\lambda|i)/(mS_m)$. The splitting identity $S_k(\lambda)=\lambda_iS_{k-1}(\lambda|i)+S_k(\lambda|i)$, applied with $k=m$ and $k=m+1$, gives
\[
 Q=\frac f m S-\frac{(m+1)f}{mS_m}S_{m+1}.
\]
For $m=n$, set $S_{n+1}=0$. If $S_{m+1}\leq0$, the second term is nonnegative. Otherwise $\lambda\in\Gamma_{m+1}$ and Maclaurin gives $S_{m+1}\leq\binom n{m+1}f^{m+1}$. Since $S_m=\binom nm f^m$, the subtracted term is at most $(n-m)f^2/m$. This proves the first assertion. Summing the derivatives gives $\mathcal F=f\sigma_{m-1}/\sigma_m\geq1$, again by Maclaurin.
\end{proof}

When $m<n$, admissible forms need not be semipositive. Positivity of a
mixed Hessian measure therefore does not directly bound every coefficient
of its factors. For degrees below $m$, deleted-coordinate positivity
supplies the required absolute bound: if $S_{j;I}(\lambda)=S_j(\lambda|I)$
denotes $S_j$ with the coordinates indexed by $I$ omitted, then
$S_{j;I}(\lambda)>0$ for $\lambda\in\Gamma_m$ and $j+|I|\leq m$, by
polarized G\aa rding positivity. By continuity, $S_{j;I}\geq0$ on
$\overline{\Gamma_m}$ in the same range.

\begin{lemma}\label{hessian:absolute-mixed-forms}
If $0\leq r<m$, $\alpha_1,\ldots,\alpha_r\in\Gamma_m(\omega)$,
and $\Xi$ is a smooth form of bidegree $(n-r,n-r)$, then
\[
 |\alpha_1\wedge\cdots\wedge\alpha_r\wedge\Xi|
 \leq C_{n,m,\Xi}(\alpha_1+\cdots+\alpha_r)^r\wedge\omega^{n-r}.
\]
For $r=0$ the right-hand side means $C_{n,m,\Xi}\omega^n$. The same
bound holds for $\alpha_1,\ldots,\alpha_r\in\overline{\Gamma_m(\omega)}$.
\end{lemma}
\begin{proof}
First take a single form $\alpha$, diagonal in an $\omega$-unitary frame,
with ordered eigenvalues $\lambda\in\Gamma_m$.
For $1\leq r<m$, repeated splitting and deleted-coordinate positivity
give $S_r(\lambda)\geq\lambda_1\cdots\lambda_r$.
The sum of the last $n-r$ eigenvalues is $S_{1;I}(\lambda)>0$, with
$I=\{1,\ldots,r\}$, so each negative
eigenvalue has modulus at most $(n-r)\lambda_{r+1}$.
Thus every product of $r$ distinct eigenvalues has modulus at most
\[
 (n-r)^r\lambda_1\cdots\lambda_r\leq(n-r)^r S_r(\lambda).
\]
Indeed, replace its negative factors by this bound and its positive
factors by the corresponding largest eigenvalues. In the diagonal frame
only complementary diagonal coefficients of $\Xi$ contribute to
$\alpha^r\wedge\Xi$. This proves
$|\alpha^r\wedge\Xi|\leq C_{n,m,\Xi}\alpha^r\wedge\omega^{n-r}$.

Apply this bound to the polynomial
$(\alpha_1+t_2\alpha_2+\cdots+t_r\alpha_r)^r\wedge\Xi$,
with $0\leq t_i\leq1$. Since
$\beta=\alpha_1+t_2\alpha_2+\cdots+t_r\alpha_r$ lies in
$\Gamma_m(\omega)$, its absolute value is at most
$C_{n,m,\Xi}\beta^r\wedge\omega^{n-r}$. Expanding
$(\alpha_1+\cdots+\alpha_r)^r=(\beta+\sum_{i\geq2}(1-t_i)\alpha_i)^r$
and using polarized positivity bounds this by the stated right-hand side.
A fixed invertible interpolation matrix on this finite-dimensional
polynomial space bounds each coefficient by its supremum on the cube.
The coefficient of $t_2\cdots t_r$ is a fixed positive multiple of
$\alpha_1\wedge\cdots\wedge\alpha_r\wedge\Xi$.
This proves the mixed bound. The closed-cone case follows by adding
$\varepsilon\omega$ to each form and taking the limit.
\end{proof}

The next coefficient bound also controls a torsion form whose
diagonal coefficients can have either sign. It compares each
coefficient $S_{q;ij}$ with two deleted coordinates with the
nonnegative coefficients $S_{q;i}$ and $S_{q;j}$ of the associated
gradient pairing.

\Needspace{14\baselineskip}
\begin{lemma}\label{hessian:deleted-newton}\index{Newton inequality!deleted coordinates}
If $\lambda\in\overline{\Gamma_m}$, $1\leq q\leq\min\{m-1,n-2\}$, and $i\ne j$,
then, with $N=n-2$ and
\[
 c_{N,q}=\frac{q(N-q)}{(q+1)(N-q+1)}<1,
\]
one has
\[
 |S_{q;ij}(\lambda)|^2
 \leq\frac{S_{q;i}(\lambda)S_{q;j}(\lambda)}{1-c_{N,q}}.
\]
\end{lemma}
\begin{proof}
First prove the unnormalized Newton inequality for an arbitrary real
list $\mu$ of length $N$. For $1\leq q<N$, put
$R(x)=\prod_{a=1}^N(x+\mu_a)$ and
$G=R^{(N-q-1)}$. Repeated Rolle's theorem makes $G$ a real-rooted
polynomial of degree $d=q+1$. Away from its roots $r_a$, logarithmic
differentiation and Cauchy--Schwarz give
\[
 G'^2-GG''=G^2\sum_{a=1}^d\frac1{(x-r_a)^2}
 \geq \frac1dG^2\left(\sum_{a=1}^d\frac1{x-r_a}\right)^2
 =\frac1dG'^2.
\]
Continuity extends this polynomial inequality to every real $x$.
At zero,
\[
 G=(N-q-1)!S_{q+1}(\mu),\quad
 G'=(N-q)!S_q(\mu),\quad
 G''=(N-q+1)!S_{q-1}(\mu).
\]
Substituting these values in $(1-1/d)G'^2\geq GG''$ gives
$S_{q-1}(\mu)S_{q+1}(\mu)\leq c_{N,q}S_q(\mu)^2$.
No sign restriction on $\mu$ or on these coefficients was used.
For $q=N$, $S_{N+1}(\mu)=0$, so the same inequality holds
with $c_{N,N}=0$.

Apply this to the list with $i,j$ deleted. Put
$A=S_{q-1;ij}$, $B=S_{q;ij}$, $C=S_{q+1;ij}$ and
$D_i=S_{q;i}$, $D_j=S_{q;j}$. Deleted-coordinate positivity gives
$A,D_i,D_j\geq0$: the required degrees are $q+1\leq m$.
Splitting in the two omitted coordinates gives the exact identity
\[
 D_iD_j=A S_{q+1}(\lambda)+B^2-AC.
\]
Since $S_{q+1}(\lambda)\geq0$ and $AC\leq c_{N,q}B^2$, this is
at least $(1-c_{N,q})B^2$.
\end{proof}

A smooth maximum preserves the cone condition and can be chosen
monotone in its smoothing scale as well as in its two arguments.

\Needspace{16\baselineskip}
\begin{lemma}\label{hessian:regularized-maximum}\index{Regularization!Hessian potentials}
If smooth uniformly bounded admissible $u_j\downarrow u$ and
$v_j\downarrow v$ are given, there are smooth uniformly bounded
admissible $w_j\downarrow\max\{u,v\}$. For any chosen
$\epsilon_j>0$ decreasing to zero, they can be chosen to satisfy
$w_j=u_j$ on $\{u_j-v_j>\epsilon_j\}$ and
$w_j=v_j$ on $\{v_j-u_j>\epsilon_j\}$.
The pointwise maximum of two bounded admissible functions is bounded
admissible. For smooth functions $f,g$ on an open subset of $X$ with
$\ddc f,\ddc g\in\overline{\Gamma_m(\omega)}$, the regularized
maximum $M_\epsilon(f,g)$ constructed below satisfies the same
unshifted cone condition.
\end{lemma}
\begin{proof}
Choose a nonnegative even smooth probability density $\zeta$
supported in $[-1,1]$, and set
\[
 r_\epsilon(t)=\int |t-\epsilon s|\zeta(s)\,\mathrm ds,
 \qquad M_\epsilon(x,y)=\tfrac12(x+y+r_\epsilon(x-y)).
\]
The function $r_\epsilon$ is smooth and convex,
$|r_\epsilon'|\leq1$, and $r_\epsilon(t)=|t|$ for
$|t|\geq\epsilon$. Jensen's inequality and the triangle inequality
give
\[
 \max\{x,y\}\leq M_\epsilon(x,y)
                     \leq\max\{x,y\}+\epsilon/2.
\]
Its partial derivatives in $x,y$ are nonnegative. For fixed $t$,
$r_\epsilon(t)$ is an even convex function of $\epsilon$, since
$\zeta$ is even. It is therefore nondecreasing for $\epsilon\geq0$.
Consequently $w_j=M_{\epsilon_j}(u_j,v_j)$ decreases pointwise to
$\max\{u,v\}$ and has the asserted equalities and bounds.

Put $h_j=u_j-v_j$ and
$a_j=(1+r_{\epsilon_j}'(h_j))/2$. The chain rule gives
\[
 \omega+\ddc w_j
 =a_j(\omega+\ddc u_j)+(1-a_j)(\omega+\ddc v_j)
       +\tfrac12r_{\epsilon_j}''(h_j)
                    \mathrm dh_j\wedge\dc h_j.
\]
Here $0\leq a_j\leq1$, and the last form is semipositive.
Convexity of the cone and its stability under semipositive addition
prove admissibility. For smooth local functions satisfying the unshifted
cone condition, omit $\omega$ from this chain rule; the same argument
proves preservation of that condition.

Finally, let $u,v$ be bounded admissible. Their pointwise maximum is
bounded and upper semicontinuous, and an upper test for it at a point
is an upper test for whichever of $u,v$ attains the maximum there.
Hence $\max\{u,v\}$ satisfies \eqref{eq:hessian-bounded-upper-tests};
in particular, it is its own canonical representative.
\end{proof}

\Needspace{11\baselineskip}
\begin{lemma}\label{hessian:canonical-order}\index{Admissible potential!pointwise order}
If bounded admissible $u,v$ satisfy $u\leq v$ almost everywhere
with respect to $\omega^n$, then $u\leq v$ everywhere.
\end{lemma}
\begin{proof}
By \cref{hessian:regularized-maximum}, $w=\max\{u,v\}$ is bounded
admissible. Since $w=v$ almost everywhere, $w$ and $v$ define the same
distribution. As bounded admissible functions, both are canonical
representatives, and a canonical representative is determined by its
distribution. Thus $w=v$ everywhere, that is, $u\leq v$.
\end{proof}

\section{Liouville and gradient estimates}
\label{sec:hessian-liouville}\index{Liouville theorem!complex Hessian cone}

The gradient estimate comes from rescaling a sequence at points of large
gradient. The background then disappears, leaving a maximal function for
a constant Euclidean cone. We first establish the Liouville theorem that
rules out such a nonconstant limit; its proof is that of
\cref{thm:stw-cone-liouville} for another cone.

For this Euclidean argument allow $1\leq m\leq n$. Write
$H(v)=(v_{i\bar j})$ for the complex Hessian. A continuous real function
$v$ on a domain in $\mathbb C^n$ is admissible for $C_{n,m}$ if
\[
 \operatorname{tr}(PH(v))\geq0
 \quad\text{in distributions for every }P\in C_{n,m}^*,
\]
where $C_{n,m}^*=\{P:\operatorname{tr}(PA)\geq0\text{ for all }A\in C_{n,m}\}$
is the dual cone of Hermitian matrices.
For smooth functions this says $H(v)\in C_{n,m}$. A \emph{lower test}
for $v$ at $z_0$ is a smooth function $\psi$ near $z_0$ with
$\psi\leq v$ near $z_0$ and $\psi(z_0)=v(z_0)$. Call $v$
\emph{maximal} if $H(\psi)(z_0)$ lies outside the interior of $C_{n,m}$
for every lower test $\psi$ at every point $z_0$. Admissibility implies subharmonicity and is
preserved by nonnegative convolution. The cone polynomial $S_m$ increases
under nonnegative matrix addition in the interior.

\begin{lemma}\label{hessian:hessian-cone-stability}
Local uniform limits of admissible maximal functions are admissible and maximal. Suppose $v$ is admissible and maximal on a bounded open set $D$, and $h$ is smooth and admissible there. If both are continuous on $\overline D$ and $h\leq v$ on $\partial D$, then $h\leq v$ on $D$.
\end{lemma}
\begin{proof}
The proof of \cref{lem:cone-maximal-stability-comparison} applies with $C_{n,m}$ in place of the cone $\{P_n(H)\geq0\}$. It uses only that distributional admissibility passes to local uniform limits, that the interior of the cone is open, and that $H+\delta I$ lies in the interior for $H$ in the cone and $\delta>0$.
\end{proof}

\begin{theorem}\label{hessian:hessian-cone-liouville}
Every bounded globally Lipschitz admissible function on $\mathbb C^n$ that is maximal for $C_{n,m}$ is constant.
\end{theorem}
\begin{proof}
For $m=1$, the dual cone consists of the nonnegative multiples of $I$. Distributional subharmonicity and the lower-test condition then say that the function is a viscosity solution of the Laplace equation. The elementary equivalence with harmonicity and bounded harmonic Liouville give the conclusion. For $m=n$, admissibility is distributional plurisubharmonicity. The continuous representative is plurisubharmonic, and bounded plurisubharmonic Liouville applies.
For a fixed $m\geq2$, we argue by induction on $n\geq m$, the case $n=m$ being the plurisubharmonic one.

For a Hermitian $(n-1)$-by-$(n-1)$ matrix $A$ and $m<n$,
\[
 S_j(A\oplus0)=S_j(A),\qquad 1\leq j\leq m.
\]
Consequently $A\oplus0$ belongs to $C_{n,m}$, respectively its interior, exactly when $A$ belongs to $C_{n-1,m}$, respectively its interior. For the closed cones, the nonnegative inequalities characterize the closure: the expansion of $S_j(\lambda+\delta\mathbf1)$ has nonnegative coefficients and a strictly positive $\delta^j$ term. If an admissible function is independent of $z_n$, convolving it and using the block identity shows that the function in the remaining variables is admissible. Extending its lower tests independently of $z_n$ shows that maximality also passes to dimension $n-1$.

With this reduction, the proof of \cref{thm:stw-cone-liouville} applies word for word, with $C_{n,m}$ in place of the cone $\{P_n(H)\geq0\}$, the induction starting at $n=m$, and \cref{hessian:hessian-cone-stability} in place of \cref{lem:cone-maximal-stability-comparison}. Apart from these two inputs, that proof uses only properties shared by $C_{n,m}$: the cone is closed, convex and stable under adding nonnegative matrices; admissibility and maximality are preserved by translations, unitary changes of coordinates and $v\mapsto av+b$ with $a>0$; and admissibility implies subharmonicity and is preserved by nonnegative convolution. For instance, for $v\geq0$ the convolution at a fixed radius of $2v_\epsilon H(v_\epsilon)$ lies in $C_{n,m}$ by convexity.
\end{proof}

A trace bound quadratic in the gradient now bounds the gradient. A
bound linear in the gradient would suffice by
\cref{lem:fixed-trace-gradient-interpolation}, since forms in
$\Gamma_m(\omega)$ have positive trace; the quadratic bound keeps its
form under rescaling, and \cref{hessian:hessian-cone-liouville} excludes
the rescaled limits.

\begin{proposition}\label{hessian:hessian-trace-to-gradient}
Let $\theta$ be a smooth real $(1,1)$-form on $X$, let $2\leq m\leq n$, and let $M,A,B>0$.
There is a uniform bound for $\sup_X|\partial u|_\omega$ over all smooth real $u$ such that $\theta_u=\theta+\ddc u\in\Gamma_m(\omega)$ and
\[
 |u|\leq M,\qquad
 \frac{\theta_u^m\wedge\omega^{n-m}}{\omega^n}\leq B,
 \qquad
 \operatorname{tr}_\omega\theta_u
       \leq A\bigl(1+\sup_X|\partial u|_\omega^2\bigr).
\]
\end{proposition}
\begin{proof}
In local holomorphic coordinates write $\omega=\mathrm i g_{i\bar j}\,\mathrm dz_i\wedge\mathrm d\bar z_j$ and $\theta_u=\mathrm i(b_{i\bar j}+\kappa u_{i\bar j})\,\mathrm dz_i\wedge\mathrm d\bar z_j$, where $\kappa=1/(2\pi)$ comes from the normalization of $\ddc$.
Since $m\geq2$, every eigenvalue of a form in $\Gamma_m(\omega)$ has modulus at most its trace, by the proof of \cref{hessian:weighted-cone-trace}. Thus the assumed trace estimate bounds every eigenvalue of $\theta_u$ relative to $\omega$ in absolute value by $A(1+\sup|\partial u|^2)$.

Suppose, contrary to the claim, that a sequence $u_j$ has maxima $R_j=\sup|\partial u_j|_\omega\to\infty$, attained at $x_j$. Passing to a subsequence, let $x_j\to x_\infty$, and work in a fixed chart, changing coordinates linearly so that $g(x_\infty)=I$. If $y_j$ is the coordinate of $x_j$, set
\begin{align*}
 v_j(z)&=u_j(y_j+z/R_j),\qquad g_j(z)=g(y_j+z/R_j),\\
 E_j(z)&=R_j^{-2}b(y_j+z/R_j)+\kappa H(v_j)(z).
\end{align*}
These are defined on balls with radii tending to infinity. On each fixed ball $g_j\to I$, the smooth background term tends to zero, and $E_j$ is uniformly bounded by the eigenvalue estimate. In particular $H(v_j)$ and its Euclidean trace are uniformly bounded.
Interior estimates for the fixed Euclidean Laplacian give local uniform $W^{2,q}$ bounds for every finite $q$. A diagonal subsequence converges in $C^1$ on compact subsets to a bounded globally Lipschitz $v$, with $|\partial v(0)|=1$.

The Hermitian matrices $g_j^{-1/2}E_jg_j^{-1/2}$ belong to the interior of $C_{n,m}$ and differ from $\kappa H(v_j)$ by a term tending uniformly to zero on fixed balls. Pairing with any fixed matrix in the dual cone and passing in distributions proves that $v$ is admissible.
By homogeneity of $S_m$, the rescaled Hessian densities satisfy
\[
 \frac{S_m(g_j^{-1/2}E_jg_j^{-1/2})}{\binom nm}
   =R_j^{-2m}\frac{\theta_{u_j}^m\wedge\omega^{n-m}}{\omega^n}
      (y_j+z/R_j)\leq BR_j^{-2m}.
\]
We prove maximality with lower tests. If a smooth lower test $\psi$ at $z_0$ had $H(\psi)(z_0)$ in the interior, choose a small $\delta>0$ so that $H(\psi)-\delta I$ remains there near $z_0$. Strictifying the contact by subtracting $\delta|z-z_0|^2$ gives interior minimum points $z_j\to z_0$ for $v_j$ minus this test. At those points
\[
 E_j\geq R_j^{-2}b+\kappa(H(\psi)-\delta I).
\]
The matrix on the right, measured with $g_j$, is uniformly inside the cone. Monotonicity of $S_m$ under nonnegative addition gives a fixed positive lower bound for the left side of the preceding equation. This contradicts $BR_j^{-2m}\to0$.
Hence $v$ is maximal. \Cref{hessian:hessian-cone-liouville} makes it constant, contradicting $|\partial v(0)|=1$.
\end{proof}

\section{A trace estimate with a prescribed potential}
\label{sec:hessian-trace}\index{Hessian equation!trace estimate}

In a coordinate chart write
$\omega=(\mathrm i/(2\pi))g_{i\bar j}\,\mathrm dz_i\wedge\mathrm d\bar z_j$.
Covariant derivatives use the Chern connection of $g$. For a smooth
strictly admissible $u$, put
\[
 H_{i\bar j}=g_{i\bar j}+u_{i\bar j},\qquad W=g^{-1}H,
 \qquad L=\sum_iF_i\nabla_{i\bar i}
\]
in a unitary frame diagonalizing $H$. Use $f,\mathcal F,S,Q$ from
\cref{sec:hessian-cone}. Since $m\geq2$, one has
$|\lambda_i|<S\leq n\lambda_1$.

We shall estimate equations whose density involves another admissible
potential $v$. In an approximation argument the derivatives of $v$ may
diverge. The logarithmic weight below absorbs these derivatives using
only uniform closeness to one fixed smooth comparison function.

\begin{lemma}\label{hessian:hessian-gradient-absorption}
For $\epsilon>0$, let $\chi(t)=-(1/12)\log(2\epsilon-t)$ on $|t|\leq\epsilon$. Suppose $q=\chi'(t)\geq8$. For any $a,b\in\mathbb C$,
\[
 \chi''(t)|a-b|^2-(9q^2+2q)|a|^2-q|a|
       \geq q^2|a|^2-1-84q^2|b|^2.
\]
In particular, let $v,\widetilde\varphi$ be smooth real functions with $|v-\widetilde\varphi|\leq\epsilon$, and take derivatives in an $\omega$-unitary frame. For a positive diagonal linearized operator with coefficients $B_p$, set $t=v-\widetilde\varphi$, $a=v_p$ and $b=\widetilde\varphi_p$, multiply by $B_p$ and sum. The result is at least $q^2\sum_pB_p|v_p|^2$, up to an error bounded by
$(1+84q^2\sup|\partial\widetilde\varphi|_\omega^2)\sum_pB_p$.
The condition $q\geq8$ holds throughout this interval if $\epsilon\leq1/288$.
\end{lemma}
\begin{proof}
Young's inequality gives $|a-b|^2\geq(7/8)|a|^2-7|b|^2$. Since $\chi''=12q^2$, the coefficient of $|a|^2$ after subtracting $9q^2$ is $3q^2/2$. Moreover, $2q\leq q^2/4$ and $q|a|\leq q^2|a|^2/4+1$. The remaining coefficient is at least $q^2$, and the error involving $b$ is $84q^2|b|^2$. Finally $q\geq1/(36\epsilon)$, which proves the last assertion. The weighted version follows by summing the scalar inequality with positive coefficients.
\end{proof}

\Needspace{18\baselineskip}
\begin{lemma}\label{hessian:fixed-comparison-absorption}
Suppose $f\geq f_0>0$, $C_H,C_G\geq0$, and $q$ satisfies
\[
 q\geq8,\qquad q\geq\frac{2nm!C_H}{f_0},\qquad
 q^2\geq\frac{4m!C_G}{f_0}.
\]
For $0\leq r\leq1/2$, $a,b\in\mathbb C^n$ and $M\in C_{n,m}$, put $W_M=\sum_iF_iM_{i\bar i}$ and
\begin{align*}
 E={}&q(W_M-\mathcal F)-\frac{C_H}{\lambda_1}|M_{1\bar1}|
             -\frac{C_G}{S}\sum_i|a_i|^2\\
 &+\sum_iF_i\bigl(12q^2|a_i-b_i|^2
             -(9q^2+2q+rq^2)|a_i|^2-q|a_i|\bigr).
\end{align*}
Then
\[
 E\geq\frac q2W_M+\frac{q^2}{4}\sum_iF_i|a_i|^2
       -(q+1+84q^2\max_i|b_i|^2)\mathcal F.
\]
\end{lemma}
\begin{proof}
The proof of \cref{hessian:hessian-gradient-absorption} uses only
$q\geq8$ and $\chi''=12q^2$. It therefore gives
\[
 12q^2|a_i-b_i|^2-(9q^2+2q)|a_i|^2-q|a_i|
 \geq q^2|a_i|^2-1-84q^2|b_i|^2.
\]
Subtracting $rq^2|a_i|^2$ leaves at least $q^2|a_i|^2/2$. Since $S\leq n\lambda_1$, \cref{hessian:weighted-cone-trace} and the hypothesis on $q$ imply
\[
 \frac{C_H}{\lambda_1}|M_{1\bar1}|
 \leq\frac{nC_H}{S}\operatorname{tr}M\leq\frac q2 W_M.
\]
The lower bound for each $F_i$ similarly gives
\[
 \frac{C_G}{S}\sum_i|a_i|^2
             \leq\frac{q^2}{4}\sum_iF_i|a_i|^2.
\]
Combining these inequalities yields the assertion. The error depends on the fixed comparison vector $b$ and is not discarded.
\end{proof}

The two cases for the smallest eigenvalue
will reduce to the following numerical estimate.

\begin{proposition}\label{hessian:conditional-maximal-eigenvalue-bound}
Fix $a>0$, $b,d,e\geq0$ and $0<\delta<1$. Let $K\geq1$, and let $\lambda\in\Gamma_m$ be ordered decreasingly, with $0<f_0\leq f\leq f_1$. Suppose the following inequality holds according to the sign of $\lambda_n+\delta\lambda_1$:
\begin{enumerate}
\item If $\lambda_n<-\delta\lambda_1$, then $0\geq aQ/K-bK\mathcal F-d\lambda_1\mathcal F-e$.
\item If $\lambda_n\geq-\delta\lambda_1$, then $0\geq aQ/K-bKF_1-e$.
\end{enumerate}
Then $S\leq CK$ and $\lambda_1\leq CK$, with $C$ depending only on the fixed displayed data and $n,m$.
\end{proposition}
\begin{proof}
In the first case, symmetry and concavity of $F$ imply $F_n\geq F_i$ for every $i$, so $Q\geq\delta^2\lambda_1^2\mathcal F/n$. To see the derivative ordering, apply monotonicity of the gradient of a concave function to $\lambda$ and the vector obtained by exchanging its $i$th and $n$th coordinates. Divide the assumed inequality by $K\mathcal F$, use $\mathcal F\geq1$, and set $t=\lambda_1/K$. It yields $a\delta^2t^2/n\leq b+dt+e$, which bounds $t$. Since $S\leq n\lambda_1$, the trace bound follows.
In the second case, since $Q\geq F_1\lambda_1^2$, either $\lambda_1^2\leq2bK^2/a$, which gives the same trace bound directly, or $aQ/(2K)\geq bKF_1$. The assumed inequality then gives $Q\leq2eK/a$. \Cref{hessian:quadratic-coercivity} bounds $S$ by a multiple of $K$. Finally the $\Gamma_2$ identity gives $\lambda_1\leq(\sum_i\lambda_i^2)^{1/2}<S$.
\end{proof}

The geometric computation uses a lower support for the largest
eigenvalue, as in the second-order argument for the Gauduchon equation
in \cref{sec:stw-second-order-proof}.
The perturbation affects only this support; derivatives of $F$ remain
evaluated at the original matrix $W$.

\Needspace{15\baselineskip}
\begin{lemma}\label{hessian:hermitian-eigenvalue-inequality}
At a point with $\lambda_1>1$, the largest eigenvalue has a smooth local lower support $\widetilde\lambda$, equal to $\lambda_1$ at the point, for which, at that point,
\[
 L\log\widetilde\lambda\geq
 \frac{f_{1\bar1}}{\lambda_1}+N
 -\sum_iF_i|Z_i|^2-C\mathcal F
 -\frac{C}{\lambda_1^2}\sum_iF_i|H_{1\bar1;i}|,
\]
where
\[
 N=-\frac{D^2F_W(\nabla_1W,\nabla_{\bar1}W)}{\lambda_1},
 \qquad Z_i=\frac{H_{i\bar1;1}}{\lambda_1}.
\]
Here $C$ depends only on the fixed metric. Derivatives of $F$ are evaluated at $W$.
\end{lemma}
\begin{proof}
Choose a smooth unitary frame parallel for the Chern connection at the point and diagonalizing $H$ there. The positive semidefinite tensor $B$, diagonal in this frame with entries $0,\beta,\ldots,\beta$, satisfies $\nabla B=0$ at the point and $|\nabla^2B|\leq C\beta$. For $0<\beta<1$, the largest eigenvalue of $W-B$ is simple, lies below that of $W$, and agrees with it at the point. This gives the required smooth support $\widetilde\lambda$. The perturbation enters only through $\widetilde\lambda$; the equation still uses $F(W)$.
Differentiating the normalized eigenvector equation gives
\[
 \widetilde\lambda_{;k}=H_{1\bar1;k},\qquad
 \widetilde\lambda_{;k\bar k}\geq H_{1\bar1;k\bar k}
 +\sum_{p>1}\frac{|H_{p\bar1;k}|^2+|H_{1\bar p;k}|^2}
                         {\lambda_1-\lambda_p+\beta}-C\beta.
\]
The denominators are at most $(n+1)\lambda_1$. Let $U=\partial\bar\partial u$, and take $T^p_{ki}=\Gamma^p_{ki}-\Gamma^p_{ik}$. The first commutations are
\[
 U_{i\bar j;k}=U_{k\bar j;i}-T^p_{ki}U_{p\bar j},\qquad
 U_{i\bar j;\bar k}=U_{i\bar k;\bar j}
                          -\overline{T^p_{kj}}U_{i\bar p}.
\]
Apply the first identity, commute one mixed pair of covariant derivatives, and apply the conjugate identity. The result has the form
\[
 U_{1\bar1;k\bar k}=U_{k\bar k;1\bar1}
      -2\operatorname{Re}\sum_p U_{k\bar p;1}\overline{T^p_{k1}}+E_k,
 \qquad |E_k|\leq C|U|.
\]
Indeed the mixed commutator contracts curvature with $U$, differentiating the first torsion coefficient contracts $\nabla T$ with $U$, and the second use of these first-order identities produces contractions of two torsion coefficients with $U$. Thus all the omitted terms satisfy the displayed bound. Since $g$ is parallel, the same derivative relations apply to $H$ apart from the explicit $g$ terms in the first commutation. At the diagonal point $|U|\leq C\lambda_1$.
For $p>1$, rewrite $U_{k\bar p;1}=U_{1\bar p;k}+T^j_{k1}U_{j\bar p}$. Young's inequality absorbs its third derivative into the positive quotient $|H_{1\bar p;k}|^2/(\lambda_1-\lambda_p+\beta)$, at the cost of $C\lambda_1$. Consequently
\[
 \widetilde\lambda_{;k\bar k}\geq
 H_{k\bar k;1\bar1}
 -2\operatorname{Re}(H_{k\bar1;1}\overline{T^1_{k1}})-C\lambda_1.
\]
Differentiating $F(W)=f$ twice in the first frame direction gives $\sum_kF_kH_{k\bar k;1\bar1}=f_{1\bar1}+\lambda_1N$. Sum the last inequality with weights $F_k$, and subtract $\sum_kF_k|\widetilde\lambda_{;k}|^2/\lambda_1^2$ to pass to $L\log\widetilde\lambda$. Write $D_k=H_{1\bar1;k}$ and $t_k=T^1_{k1}$. The first commutation, using $U_{1\bar1}=\lambda_1-1$, gives
\[
 D_k=\lambda_1Z_k-(\lambda_1-1)t_k.
\]
The negative logarithmic square and the remaining torsion term therefore satisfy
\[
 -\frac{|D_k|^2}{\lambda_1^2}-2\operatorname{Re}(Z_k\bar t_k)
 \geq-|Z_k|^2-C-C\lambda_1^{-2}|D_k|.
\]
For example, expanding the left side leaves $-|Z_k|^2-2\lambda_1^{-1}\operatorname{Re}(Z_k\bar t_k)-(1-\lambda_1^{-1})^2|t_k|^2$, and $\lambda_1|Z_k|\leq|D_k|+C\lambda_1$. This proves the claim, with $C$ independent of $\beta\in(0,1)$.
\end{proof}

\begin{lemma}\label{hessian:hermitian-gradient-commutation}
Put $s=|\partial u|_g^2$ and $K=1+\sup_Xs$. For every $\eta>0$ one has
\[
 Ls\geq\sum_{i,k}F_i(|u_{ki}|^2+|u_{k\bar i}|^2)
 +2\operatorname{Re}\sum_k f_k u_{\bar k}-\eta Q-C_\eta K\mathcal F.
\]
The constant depends only on the metric and $\eta$.
\end{lemma}
\begin{proof}
Differentiate the squared norm using metric compatibility. Its quadratic part is the displayed sum. Commuting the remaining third derivatives gives
\[
 u_{ki\bar i}=u_{i\bar i;k}-T^p_{ik}u_{p\bar i}
                  +R^p_{i\bar i k}u_p,
 \qquad
 u_{\bar k i\bar i}=u_{i\bar i;\bar k}
                  -\overline{T^p_{ik}}u_{i\bar p}.
\]
The first differentiated equation is $\sum_iF_i u_{i\bar i;k}=f_k$, because the background $g$ is parallel. At the diagonal point the torsion errors are bounded by $C\sqrt K\sum_iF_i(|\lambda_i|+1)$, and the curvature errors by $CK\mathcal F$. Use
\[
 C\sqrt K F_i|\lambda_i|
       \leq\eta' F_i\lambda_i^2+C_{\eta'}KF_i
\]
with a sufficiently small $\eta'$ for the finite sum. This gives the asserted bound.
\end{proof}

\Needspace{19\baselineskip}
\begin{proposition}\label{hessian:prescribed-potential-trace}
Fix $M,B\geq1$ and a smooth real $h(x,t)$ on a neighborhood of
$X\times[-B-1,B+1]$. There is $0<\varepsilon\leq1$, depending only on
$M,B,h,\omega,n,m$, with the following property. For every smooth real
$\varphi$ with $\|\varphi\|_\infty\leq B$, there is $C<\infty$, depending
also on its $C^2$ norm, such that if smooth real $u,v$ satisfy
\begin{gather*}
 W\in\Gamma_m,\qquad I+v_{i\bar j}\in\overline{\Gamma_m},\\
 \|u\|_\infty\leq M,\qquad\|v-\varphi\|_\infty\leq\varepsilon,\\
 F(W)=\exp\bigl((u+h(x,v))/m\bigr),
\end{gather*}
then
\[
 \operatorname{tr}_\omega(\omega+\ddc u)
                  \leq C\bigl(1+\sup_X|\partial u|_g^2\bigr).
\]
The cone condition on $v$ is understood in a unitary frame. The tolerance $\varepsilon$ uses no derivatives of $\varphi$; the constant $C$ uses no derivatives of $v$.
\end{proposition}
\begin{proof}
For $\varepsilon\leq1$ the density has fixed bounds $0<f_0\leq f\leq f_1$. Its first derivative is
\[
 f_k=\frac f m(u_k+h_k(x,v)+h_t(x,v)v_k).
\]
In \cref{hessian:hermitian-gradient-commutation}, the $u_k$ terms contribute $2fs/m\geq0$ to $2\operatorname{Re}\sum_kf_ku_{\bar k}$, so they can be dropped. At a point with $\lambda_1>1$, differentiating once more gives
\[
 f_{1\bar1}\geq-C(1+|\partial v|^2+|1+v_{1\bar1}|).
\]
Here $u_{1\bar1}=\lambda_1-1>0$, the square from differentiating the exponential is nonnegative, and all derivatives of $h$ are bounded on the fixed interval.

Set
\begin{align*}
 \Phi(s)&=-\tfrac12\log(1-s/(2K)),\qquad
 \Psi(u)=-2Au+(A\tau/2)u^2,\\
 \chi(t)&=-\tfrac1{12}\log(2\varepsilon-t),\qquad \tau=(M+1)^{-1}.
\end{align*}
Then $(4K)^{-1}\leq\Phi'\leq(2K)^{-1}$ and $\Phi''=2(\Phi')^2$. Since $|u|\leq M$ and $\tau M<1$, one has $-3A\leq\Psi'\leq-A$, and $\Psi''=A\tau$. Finally $q=\chi'$ lies between $1/(36\varepsilon)$ and $1/(12\varepsilon)$, with $\chi''=12q^2$.
Take a maximum point of $G=\log\lambda_1+\Phi(s)+\Psi(u)+\chi(v-\varphi)$. Replacing $\lambda_1$ by the smooth lower support $\widetilde\lambda$ of \cref{hessian:hermitian-eigenvalue-inequality} gives a smooth function $\widetilde G\leq G$ with a local maximum at this point. There put
\[
 a_i=v_i,\quad b_i=\varphi_i,\quad P_v=I+v_{i\bar j},\quad
 W_v=\sum_iF_i(P_v)_{i\bar i},\quad
 x_i=\Phi's_i,\quad y_i=\Psi'u_i.
\]
The first-order condition at the maximum, and the identity $D_i=\lambda_1Z_i-(\lambda_1-1)T^1_{i1}$ from the proof of \cref{hessian:hermitian-eigenvalue-inequality}, give
\[
 H_{1\bar1;i}/\lambda_1=-x_i-y_i-q(a_i-b_i),\qquad
 Z_i=-x_i-y_i-qa_i+e_i,\quad
 e_i=qb_i+(1-\lambda_1^{-1})T^1_{i1}.
\]
Thus $|e_i|\leq E_0$, with $E_0$ independent of $A$ and of derivatives of $v$.

We first bound the term $\Phi'Ls$ of $LG$. \Cref{hessian:hermitian-gradient-commutation} and the identity $u_{i\bar j}=(\lambda_i-1)\delta_{ij}$ give, after multiplication by $\Phi'$,
\[
 \Phi'Ls\geq\frac1{8K}\sum_{i,k}F_i(|u_{ki}|^2+|u_{k\bar i}|^2)
             +\frac{3}{64K}Q-C\mathcal F
             -\frac{C_G'}S\sum_i|a_i|^2.
\]
To justify the constants, retain half of the positive quadratic term, and bound the other half below by its diagonal terms $|u_{i\bar i}|^2$, using $(\lambda_i-1)^2\geq\lambda_i^2/2-1$. Choose $\eta$ in that lemma so that the term $\eta\Phi'Q$ costs at most $Q/(128K)$. The remaining density-gradient error is at most $C|a|/\sqrt K$ plus a fixed constant. Young's inequality bounds $C|a|/\sqrt K$ by $C_G'|a|^2/S+cS/K$. Choose $C_G'$ so large that, by \cref{hessian:quadratic-coercivity}, $cS/K$ is at most $Q/(128K)$ plus a fixed constant. This leaves $3Q/(64K)$; fixed constants are at most $C\mathcal F$, since $\mathcal F\geq1$.
The last error term of \cref{hessian:hermitian-eigenvalue-inequality}, linear in $H_{1\bar1;i}$, is estimated using the first-order condition. Its $x_i$ term costs at most half of the retained quadratic sum and $C\mathcal F$. Its $y_i$ term costs at most $\Psi''\sum_iF_i|u_i|^2/4+\mathcal F$ whenever $\lambda_1\geq L_0(A)$, where $L_0(A)$ is fixed independently of $K,v$. Its $qa_i$ term costs at most $q\sum_iF_i|a_i|$ when $\lambda_1$ exceeds a fixed metric constant, which we take at least $1$; its $qb_i$ term costs $C\mathcal F$.
If $\lambda_1$ lies below either threshold at the maximum point, the final comparison of this proof gives the required bound directly. Otherwise, with fixed $C_H,C_G,C_0$ independent of $A$, the second-order condition $L\widetilde G\leq0$ gives
\begin{align*}
 0\geq{}&N-\sum_iF_i|Z_i|^2+2\sum_iF_i|x_i|^2
       +\tfrac34\Psi''\sum_iF_i|u_i|^2+\frac{Q}{32K}
       +\Psi'(f-\mathcal F)-C_0\mathcal F\\
 &+q(W_v-\mathcal F)-\frac{C_H}{\lambda_1}|(P_v)_{1\bar1}|
       -\frac{C_G}S\sum_i|a_i|^2
       +\sum_iF_i\bigl(12q^2|a_i-b_i|^2-q|a_i|\bigr).
\end{align*}
The fixed $qL\varphi$ term is included in $C_0$. Choose $\varepsilon$ so small that the lower bound $q\geq1/(36\varepsilon)$ meets all thresholds of \cref{hessian:fixed-comparison-absorption}. These thresholds depend on $M,B,h$ and the metric, and not on derivatives of $\varphi$. This choice precedes that of $A$.

First suppose $\lambda_n<-\delta\lambda_1$. Concavity gives $N\geq0$, and
\[
 -|Z_i|^2+2|x_i|^2
     \geq-6|y_i|^2-6q^2|a_i|^2-6E_0^2.
\]
\Cref{hessian:fixed-comparison-absorption}, with $M=P_v$ and $r=0$, applies after decreasing the right side further by $(3q^2+2q)|a_i|^2$. Drop the nonnegative terms $qW_v/2$ and $q^2\sum_iF_i|a_i|^2/4$ that it provides, and the $\Psi''$ term. Since $\Psi'f\geq-3Af_1$ and $-\Psi'\mathcal F\geq0$, this gives
\[
 0\geq\frac{Q}{32K}-(54A^2+C_*)K\mathcal F-3Af_1,
\]
where $C_*$ is fixed independently of $A,\delta$.

Now suppose $\lambda_n\geq-\delta\lambda_1$, and let $I=\{i:F_i>\delta^{-1}F_1\}$. For $i\in I$, $\lambda_i<\lambda_1$ and
\[
 \frac{F_i-F_1}{\lambda_1-\lambda_i}
       \geq\frac{(1-\delta)F_i}{(1+\delta)\lambda_1}
       \geq\frac{(1-2\delta)F_i}{\lambda_1}.
\]
The spectral formula for $D^2F$, including its nonpositive diagonal part, therefore gives $N\geq(1-2\delta)\sum_{i\in I}F_i|Z_i|^2$. This uses the original eigenvalues; repeated eigenvalues do not occur in $I$.
Using $x_i=-Z_i-y_i-qa_i+e_i$ and $|z+w|^2\geq\delta|z|^2-\delta|w|^2/(1-\delta)$, one obtains
\begin{align*}
 N-\sum_{i\in I}F_i|Z_i|^2+2\sum_{i\in I}F_i|x_i|^2
 &\geq-\frac{4\delta(\Psi')^2}{1-\delta}\sum_{i\in I}F_i|u_i|^2\\
 &\quad-r q^2\sum_{i\in I}F_i|a_i|^2-rE_0^2\mathcal F,
 \qquad r=\frac{8\delta}{1-\delta}.
\end{align*}
Choose $\delta\leq\min\{1/17,\tau/(72A+\tau)\}$. Then the first coefficient is at most $\Psi''/2$, and $r\leq1/2$. On the complementary indices the preceding elementary bound gives
\[
 -\sum_{i\notin I}F_i|Z_i|^2+2\sum_{i\notin I}F_i|x_i|^2
 \geq-\frac{54A^2+6nE_0^2}{\delta}KF_1
       -6q^2\sum_{i\notin I}F_i|a_i|^2.
\]
Apply \cref{hessian:fixed-comparison-absorption} with $M=P_v$ and penalty $r$. To match its input, add the negative terms $-(9q^2+2q)|a_i|^2$ for $i\in I$ and $-(3q^2+2q+rq^2)|a_i|^2$ for $i\notin I$; this only decreases the right side. Its fixed-comparison error and the error $rE_0^2\mathcal F$ from $I$ have constants bounded independently of $A,\delta$. Enlarge $C_*$ to include them, and drop the nonnegative terms provided by the lemma and the remaining $\Psi''$ term. We obtain
\[
 0\geq\frac{Q}{32K}-\frac{54A^2+6nE_0^2}{\delta}KF_1
                +(A-C_*)\mathcal F-3Af_1.
\]
Take $A\geq\max\{1,C_*+1\}$, then fix $\delta=\min\{1/17,\tau/(72A+\tau)\}$ and the threshold $L_0(A)$. Both case inequalities now satisfy the hypotheses of \cref{hessian:conditional-maximal-eigenvalue-bound}, with fixed data. Hence $\lambda_1\leq CK$ at the maximum point.
Finally $\Phi,\Psi,\chi$ have oscillations bounded independently of $K$ and of derivatives of $v$. Comparing $G$ at any other point with its maximum gives the same bound for $\lambda_1$ everywhere. Since $S\leq n\lambda_1$, this proves the trace estimate.
\end{proof}

\begin{corollary}\label{hessian:prescribed-potential-complex-hessian}
Under the hypotheses of \cref{hessian:prescribed-potential-trace}, one has
\[
 \sup_X|\partial u|_g+\sup_X|\partial\bar\partial u|_g\leq C,
\]
with a constant depending on the same fixed data and no derivatives of $v$.
\end{corollary}
\begin{proof}
The Hessian density is $F(W)^m\leq f_1^m$. Apply \cref{hessian:hessian-trace-to-gradient} to the fixed background $\theta=\omega$, with the trace estimate of \cref{hessian:prescribed-potential-trace}. It gives the uniform gradient bound. The trace is then uniformly bounded, and the $\Gamma_2$ identity bounds the absolute value of every eigenvalue of $W$ by its trace. Subtracting the identity matrix gives the complex-Hessian bound.
\end{proof}

The bound controls the complex Hessian. An estimate for all entries of
the real Hessian is not needed: the Euclidean Laplacian supplies finite
$L^p$ estimates for those derivatives.

\section{Bounded complex Hessians and elliptic extensions}
\label{sec:hessian-complex-limits}
\index{Hessian equation!uniformly elliptic extension}

Uniform complex-Hessian bounds give enough compactness both to pass the
nonlinear products to a limit and to make the Hessian root uniformly
elliptic. In coordinates use the metric matrix $g$ of the preceding
section and write $H(u)=(u_{i\bar j})$, as in \cref{sec:hessian-liouville}. Coefficient products with bounded
measurable factors will be identified with currents by convolution.

\begin{lemma}\label{hessian:complex-hessian-compactness}
Suppose $u_j$ is smooth,
\[
 \sup_j\bigl(\|u_j\|_\infty+\|\ddc u_j\|_{L^\infty}\bigr)<\infty,
 \qquad u_j\longrightarrow u\quad\text{uniformly on }X.
\]
Then $u\in W^{2,p}(X)$ for every $1<p<\infty$, its distributional complex
Hessian has bounded coefficients, and $u_j\to u$ in $C^{1,\alpha}(X)$
for every $0<\alpha<1$. The same hypotheses give uniform $W^{2,p}$ bounds
for the sequence.
\end{lemma}
\begin{proof}
On a coordinate ball the Euclidean Laplacian is
$\Delta u_j=4\sum_i (u_j)_{i\bar i}$, so it is uniformly bounded.
The interior Laplace estimate on two nested balls gives
\[
 \|u_j\|_{W^{2,p}(B')}\leq C_p
       \bigl(\|\Delta u_j\|_{L^p(B)}+\|u_j\|_{L^p(B)}\bigr).
\]
A finite cover gives the global bound. Weak $W^{2,p}$ limits coincide
with the uniform limit. The bounded complex Hessians have weak-star
limits in $L^\infty$; testing against compactly supported smooth functions
identifies them with $H(u)$. Choose $p$ with $1-2n/p>\alpha$.
The compact Sobolev embedding into $C^{1,\alpha}$ gives convergence of
subsequences. Every such limit is $u$, hence the whole sequence converges.
This estimates the full real Hessian in every finite $L^p$ space, without
asserting that it is bounded.
\end{proof}

\begin{lemma}\label{hessian:complex-hessian-products}
Let $1\leq k\leq n$. For $1\leq a\leq k$, suppose that $u_j^a$ and $u^a$ satisfy
the hypotheses of \cref{hessian:complex-hessian-compactness}. Let $\theta_a$ be fixed smooth real $(1,1)$-forms. Then
\[
 \bigwedge_{a=1}^k(\theta_a+\ddc u_j^a)
       \longrightarrow
 \bigwedge_{a=1}^k(\theta_a+\ddc u^a)
       \quad\text{as currents}.
\]
The current on the right is the product of the almost-everywhere
coefficients. No closedness or positivity of the $\theta_a$ is required.
For each smooth test form $\eta$ of bidegree $(n-k,n-k)$, the
difference of the pairings with $\eta$ is bounded by
$C_\eta\max_a\|u_j^a-u^a\|_\infty$.
\end{lemma}
\begin{proof}
Set $D_j^a=\ddc u_j^a$ and $D^a=\ddc u^a$.
On a smaller coordinate ball, convolution gives smooth $u_\epsilon^a$
whose complex Hessians are uniformly bounded and converge to $H(u^a)$
in every finite $L^p$ space. Their coefficient products consequently
converge in finite $L^p$ spaces. In particular the products of the $D^a$
are closed currents, since their smooth approximants are closed.

Expand the smooth background factors and telescope each remaining
product. Every difference is then a sum of terms of the form
\[
 \ddc(u_j^a-u^a)\wedge T_j,
\]
where $T_j$ is a product of $D_j^b$ and $D^b$ with $b\ne a$.
It has uniformly bounded coefficients and is closed.
Local convolution, followed by strong convergence of its finite
coefficient products, justifies integration by parts:
\[
 \left\langle\ddc(u_j^a-u^a)\wedge T_j,\xi\right\rangle
       =\int (u_j^a-u^a)\,T_j\wedge\ddc\xi.
\]
Here $\xi$ includes the chosen smooth background factors and the test
form, so their derivatives are retained. The right-hand side has the
asserted uniform bound. A finite partition of unity proves the claim on
$X$. This argument also proves that the coefficient product is the
distributional limit; it does not identify a separately defined bounded
Hessian operator for all nonsmooth admissible functions.
\end{proof}

\Needspace{22\baselineskip}
\begin{lemma}\label{hessian:uniform-elliptic-extension}
For $f_0>0$ and $L>0$, assume the set
\[
 K=\{A\in C_{n,m}:\operatorname{tr}A\leq L,
                         \ \sigma_m(A)\geq f_0^m\}
\]
is nonempty. It is a compact subset of the interior, and there are
$\vartheta,\Theta>0$, depending only on $n,m,f_0,L$, such that
\[
 \vartheta I\leq DF_A\leq\Theta I\qquad(A\in K).
\]
The function on all Hermitian matrices defined by
\[
 \widehat F(B)=\inf_{A\in K}
       \bigl(F(A)+\operatorname{tr}(DF_A(B-A))\bigr)
\]
is concave, agrees with $F$ on $K$, and satisfies, for $P\geq0$,
\[
 \vartheta\operatorname{tr}P\leq
       \widehat F(B+P)-\widehat F(B)
       \leq\Theta\operatorname{tr}P.
\]
\end{lemma}
\begin{proof}
In the closed $2$-cone,
\[
 \sum_i\lambda_i^2=(\operatorname{tr}A)^2-2S_2(A)
                          \leq L^2.
\]
Thus $K$ is bounded and closed. Maclaurin gives
$\sigma_j(A)\geq\sigma_m(A)^{j/m}\geq f_0^j$ for $j\leq m$,
so every matrix in $K$ is in the interior. In a unitary eigenframe,
\[
 F_i=\frac{F^{1-m}S_{m-1}(\lambda|i)}{m\binom nm},
 \qquad
 \frac{f_0}{m!L}\leq F_i\leq
                 \frac{L^{m-1}}{n f_0^{m-1}}.
\]
The lower bound is \cref{hessian:weighted-cone-trace}; the upper bound
uses $|\lambda_i|\leq L$ in the polynomial numerator and its positivity.
These bounds give possible values of $\vartheta$ and $\Theta$.
Concavity places each tangent plane above $F$ on the cone; the plane
at $B$ gives equality when $B\in K$. An infimum of affine functions is
concave. Each plane increases under $P\geq0$ by a number between the
two displayed trace bounds; taking infima preserves these bounds.
The same estimates give global Lipschitz continuity of $\widehat F$.
\end{proof}

To interpret this as a real elliptic operator, let $\mathcal C(R)$ be
the complex Hessian of the real quadratic function
$q(x)=\frac12 x^{\mathsf T}Rx$, where $x\in\mathbb R^{2n}$.
For $R\geq0$, one has $\mathcal C(R)\geq0$ and
$\operatorname{tr}\mathcal C(R)=\operatorname{tr}R/4$.
If $\mu I\leq g(x)\leq MI$, then
\[
 \mathcal G(x,R)=\widehat F
       \bigl(I+g(x)^{-1/2}\mathcal C(R)g(x)^{-1/2}\bigr)
\]
is uniformly elliptic on real symmetric matrices, with constants
$\vartheta/(4M)$ and $\Theta/(4\mu)$.
Its dependence on $x$ satisfies on compact coordinate balls
\[
 |\mathcal G(x,R)-\mathcal G(y,R)|
                      \leq C|x-y|(1+\|R\|).
\]
The ellipticity follows from the trace identity for $\mathcal C$, and the
$x$-dependence from the Lipschitz continuity of $\widehat F$ and the smooth
dependence of the positive matrix square root of the fixed metric.

We use the classical interior Evans--Krylov--Caffarelli estimate for a
concave operator on all real symmetric matrices with fixed positive
ellipticity constants and coefficient oscillation bounded by
$C|x-y|^\gamma(1+\|R\|)$, where $\gamma>0$. A H\"older right-hand side
then gives an interior $C^{2,\beta}$ estimate, for some $\beta>0$.
This estimate applies also to continuous viscosity solutions. The real
extension just constructed satisfies these conditions. Once its solution
is $C^{2,\beta}$ and the attained complex Hessians remain in a compact
subset of the open cone, the original smooth Hessian operator and
Schauder estimates give the higher derivatives.

We next identify the limit of the proper auxiliary equations that will
be used for regularity.

\begin{proposition}\label{hessian:proper-auxiliary-limit}
Suppose $v_j\to\varphi$ uniformly on $X$, $h(x,t)$ is smooth on a
neighborhood of $X\times J$ for a compact interval $J$ containing all
their values, and smooth admissible $u_j$ satisfy
\[
 (\omega+\ddc u_j)^m\wedge\omega^{n-m}
           =e^{u_j+h(x,v_j)}\omega^n.
\]
Assume $u_j\to u$ uniformly and the complex Hessians of $u_j$ are uniformly
bounded. Then $\omega+\ddc u\in\overline{\Gamma_m(\omega)}$ almost everywhere, and
its coefficient product satisfies
\[
 (\omega+\ddc u)^m\wedge\omega^{n-m}
           =e^{u+h(x,\varphi)}\omega^n.
\]
On every coordinate ball $u$ also solves the corresponding uniformly
elliptic concave extension in the viscosity sense.
\end{proposition}
\begin{proof}
\Cref{hessian:complex-hessian-compactness} gives $u\in W^{2,p}$ for every finite $p$ and a bounded
complex Hessian. The matrices
$A_j=I+g^{-1/2}H(u_j)g^{-1/2}$ converge weakly in $L^2$ to the analogous
matrix $A$ for $u$. Mazur's lemma gives convex combinations of tails
converging strongly in $L^2$, hence almost everywhere along a subsequence.
The pointwise closed cone is convex, so $A\in C_{n,m}$ almost everywhere.
\Cref{hessian:complex-hessian-products} passes the equation to the limit as currents, since its
right-hand side converges uniformly. Equality of the bounded densities
then gives the equation almost everywhere.

The right-hand sides are bounded above and away from zero. The complex
Hessian bound gives a uniform upper bound $L$ for $\operatorname{tr}A_j$.
Choose $f_0>0$ below all their normalized roots. Thus every $A_j$, and
almost every $A$, belongs to the compact set $K$ of \cref{hessian:uniform-elliptic-extension}.
Every $u_j$ is a classical solution of
\[
 \mathcal G(x,D^2u_j)=r_j(x),\qquad
               r_j(x)=e^{(u_j+h(x,v_j))/m}.
\]
For completeness, let a smooth test $q$ touch $u$ from above at an
interior point $x_0$. Replacing $q$ by $q+|x-x_0|^4$ makes the local
contact strict and does not change $D^2q(x_0)$. Local maximum points
$x_j$ of $u_j-q$ then tend to $x_0$. At these points $D^2u_j\leq D^2q$,
and ellipticity gives $\mathcal G(x_j,D^2q(x_j))\geq r_j(x_j)$. Passing to
the limit gives $\mathcal G(x_0,D^2q(x_0))\geq e^{(u+h(x,\varphi))/m}$ at
$x_0$. A test from below gives the reverse inequality. The extension is
defined on every real Hessian, so no admissibility of a lower test has
been assumed. This proves the viscosity statement.
\end{proof}

\section{Proper equations and admissible approximation}
\label{sec:hessian-proper-approximation}
\index{Hessian equation!proper exponential equation}
\index{Admissible potential!smooth approximation}

For a smooth potential write
$H_m(w)=(\omega+\ddc w)^m\wedge\omega^{n-m}$.
The positive exponential coefficient controls additive constants by the
maximum principle. The resulting smooth equations will also construct
admissible approximations of continuous potentials, before any continuous
Hessian product is defined.

\Needspace{11\baselineskip}
\begin{lemma}\label{hessian:smooth-proper-existence}
For every $a>0$ and smooth real $H$ on $X$, there is a unique smooth
strictly admissible solution of
\[
 H_m(w)=e^{aw+H}\omega^n.
\]
It satisfies
\[
 -\frac{\max_X H}{a}\leq w\leq-\frac{\min_X H}{a}.
\]
\end{lemma}
\begin{proof}
At a maximum of $w$, the admissible form $\omega+\ddc w$ is at most
$\omega$, hence $e^{aw+H}\leq1$ there. At a minimum it is at least
$\omega$, and the reverse inequality holds. These give the bounds.
For two solutions, a positive maximum of their difference would imply
$H_m(w_1)\leq H_m(w_2)$ there, contradicting the exponential factor.
This proves uniqueness.

First take $a=1$ and use the path
$H_m(w_t)=e^{w_t+tH}\omega^n$, starting with $w_0=0$.
The trace estimate of \cref{hessian:prescribed-potential-trace} applies with prescribed potential $v=0$,
fixed smooth comparison $0$, and $h(x,v)=tH(x)$; its constants are
uniform for $t\in[0,1]$. \Cref{hessian:hessian-trace-to-gradient} then gives a
uniform gradient and complex-Hessian bound. \Cref{hessian:complex-hessian-compactness} gives uniform $W^{2,p}$ and hence $C^{1,\alpha}$ bounds, for any
fixed $\alpha<1$. Consequently
$r_t=e^{(w_t+tH)/m}$ has a uniform H\"older bound.

The extension of \cref{hessian:uniform-elliptic-extension} is uniformly elliptic and concave on all real
Hessians, with coefficient oscillation bounded by
$C|x-y|(1+\|R\|)$. The analytic estimate recalled in \cref{sec:hessian-complex-limits} therefore gives uniform $C^{2,\beta}$
bounds on a finite cover. The original smooth Hessian operator is
uniformly elliptic near the attained compact cone set. Differentiating
and applying Schauder estimates gives bounds for every higher derivative.

The linearization of
$\log(H_m(w)/\omega^n)-w-tH$ is $L_w-1$, where
\[
 L_w\xi=\frac{m\ddc\xi\wedge(\omega+\ddc w)^{m-1}
                                  \wedge\omega^{n-m}}{H_m(w)}.
\]
Its principal coefficients are positive definite. The maximum principle gives zero kernel;
smooth elliptic Fredholm theory gives surjectivity and Schauder inversion.
The implicit-function theorem proves openness. The uniform estimates
prove closedness, so the path reaches $t=1$.
For general $a$, use the metric $a\omega$ and potential $\widehat w=aw$.
Both sides of its coefficient-one equation acquire the same factor
$a^n$, so it is exactly the stated equation after division by that factor.
\end{proof}

An ordinary smooth approximation need not be admissible. To obtain an
admissible one, choose a smooth obstacle above the continuous potential
and penalize the solution's distance from that obstacle.

\begin{lemma}\label{hessian:continuous-smooth-approximation}
Every continuous admissible $\varphi$ on $X$ is the uniform limit of
smooth strictly admissible functions.
\end{lemma}
\begin{proof}
Fix a smooth $h\geq\varphi$ and a constant
\[
 b>\max\left\{1,\left\|H_m(h)/\omega^n\right\|_\infty\right\}.
\]
For each positive integer $N$ solve, by the preceding lemma,
\[
 H_m(w_N)=b e^{N(w_N-h)}\omega^n.
\]
At a positive maximum of $w_N-h$, the form $\omega+\ddc h$ lies above
the strictly admissible form $\omega+\ddc w_N$, hence is admissible there.
Monotonicity would give
$b e^{N(w_N-h)}\leq H_m(h)/\omega^n<b$, a contradiction.
Thus $w_N\leq h$. The preceding zero-order estimate also gives
\[
 \|w_N\|_\infty\leq M_N:=\|h\|_\infty+\frac{\log b}{N}.
\]

Let $x_N$ maximize $\varphi-(1+N^{-1})w_N$.
The smooth positive form
$T_N=(\omega+\ddc w_N)^{m-1}\wedge\omega^{n-m}$
defines a strictly elliptic linear operator. Admissibility gives
$(\omega+\ddc\varphi)\wedge T_N\geq0$ in distributions.
Its linear maximum principle, applied to the upper test
$(1+N^{-1})w_N$ plus the maximum value, gives at $x_N$
\[
 0\leq\bigl(\omega+(1+N^{-1})\ddc w_N\bigr)\wedge T_N.
\]
This application is to a continuous distributional subsolution of a
smooth strictly elliptic operator. It does not assume a second derivative
of $\varphi$. Writing $d_N=H_m(w_N)/\omega^n$, polarized G\aa rding
positivity yields
\[
 (N+1)d_N\geq \omega\wedge T_N/\omega^n
                         \geq d_N^{(m-1)/m}.
\]
Therefore, at this maximum point,
\[
 w_N-h\geq-E_N,\qquad
 E_N=\frac{m\log(N+1)+\log b}{N}.
\]
Since $\varphi\leq h$, the maximum of
$\varphi-(1+N^{-1})w_N$ is at most $E_N+M_N/N$.
It follows on all of $X$ that
\[
 \varphi-E_N-2M_N/N\leq w_N\leq h.
\]
Given $\delta>0$, choose an ordinary smooth majorant
$\varphi\leq h\leq\varphi+\delta$, and then $N$ so large that
$E_N+2M_N/N\leq\delta$. The resulting $w_N$ is strictly admissible
and within $\delta$ of $\varphi$. Let $\delta\downarrow0$.
No monotonicity of the chosen approximants is needed.
\end{proof}

\section{Continuous products, envelopes and domination}
\label{sec:hessian-continuous-products}\index{Hessian measure!continuous potentials}

Uniform admissible approximation allows us to define the mixed Hessian
products. Their local bounds retain the derivatives of the Hermitian
background; closedness of the background is not required.

\Needspace{13\baselineskip}
\begin{proposition}\label{hessian:continuous-products}
For continuous admissible $u_1,\ldots,u_k$, $1\leq k\leq m$, the current
\[
 (\omega+\ddc u_1)\wedge\cdots\wedge(\omega+\ddc u_k)
                                          \wedge\omega^{n-m}
\]
is well defined as the limit of uniform smooth admissible approximations.
It is positive, local, and continuous when its potentials converge
uniformly. Its local masses are bounded in terms of uniform bounds on
the potentials. In particular $H_m(u)$ is a positive Radon measure.
\end{proposition}
\begin{proof}
\Cref{hessian:smooth-proper-existence,hessian:continuous-smooth-approximation} provide the smooth approximations without
using any continuous Hessian product or domination input.
For a smooth admissible $u$ put $\sigma=\omega+\ddc u$.
Since $\mathrm d\sigma=\mathrm d\omega$, the coefficients of
$\ddc(\chi\sigma^{r-1}\wedge\omega^{n-r})$ are sums of smooth test
factors paired with powers $\sigma^q$, $q\leq r-1$, where $\chi$ is a
fixed cutoff. The absolute-form lemma and integration by parts give
the Chern--Levine--Nirenberg bound by induction on $r\leq m$:
\[
 \int_K\sigma^r\wedge\omega^{n-r}
                     \leq C_{K,U,\omega}(1+\|u\|_{L^\infty(U)})^r.
\]
At the induction step write $\sigma^r=\omega\wedge\sigma^{r-1}
+\ddc u\wedge\sigma^{r-1}$, integrate the second term against $\chi$,
and move $\ddc$ to $\chi\sigma^{r-1}\omega^{n-r}$.
The resulting powers have degrees $r-1,r-2,r-3$, with negative powers absent, and the induction
applies on nested slightly larger compact sets. Mixed masses reduce to
the single potential average, by positivity of the expansion of its
power. Positive mixed currents have total variation controlled by these
trace masses.

To prove independence, telescope the difference between two smooth
mixed products whose potentials have a common uniform bound. Each term
paired with a fixed smooth complementary test is
\[
 \int\ddc(u_a-v_a)\wedge T\wedge\eta
                         =\int(u_a-v_a)\ddc(T\wedge\eta),
\]
where $T$ is a product of the remaining admissible factors and
$\omega^{n-m}$. The identity $\mathrm d(\omega+\ddc u)=\mathrm d\omega$
and its $\dc$ analogue show that derivatives of $T$ introduce only
fixed metric derivatives; every term has at most $k-1<m$ admissible
factors. The absolute mixed-form bound and the mass estimate just
proved bound the variation of $\ddc(T\wedge\eta)$ uniformly.
Consequently the pairing difference is at most
$C_\eta\max_a\|u_a-v_a\|_\infty$.
The same construction with an arbitrary complementary smooth test
also defines the signed products of fewer than $m$ factors without the
factor $\omega^{n-m}$: their variation is bounded by the absolute-form
lemma. These products justify the metric-derivative terms above.
This proves that all approximations have the same limit and that the
products are continuous in the uniform topology. Positivity passes
from the smooth mixed products. Using a test supported in an open set
on which two limiting potentials coincide proves locality by the same
estimate on a smaller neighborhood of its support.
\end{proof}

Write $H_m(u)$ for the measure with every potential equal to $u$.
Locality and uniform continuity of these products give comparison on a
contact set.

\begin{lemma}\label{hessian:contact-comparison}
If $u\leq v$ are continuous admissible functions and the remaining
$m-1$ factors are fixed continuous admissible functions, then on
$\{u=v\}$ the mixed measure with first factor $\omega+\ddc u$ is at
most the one with first factor $\omega+\ddc v$.
In particular $\mathbf1_{\{u=v\}}H_m(u)\leq
\mathbf1_{\{u=v\}}H_m(v)$.
\end{lemma}
\begin{proof}
Write $M(w)$ for the mixed measure with first potential $w$.
Positivity and locality give, for $t>0$,
\[
 M(\max(v,u+t))\geq
 \mathbf1_{\{v>u+t\}}M(v)+\mathbf1_{\{u\geq v\}}M(u).
\]
The first indicators increase to $\mathbf1_{\{v>u\}}$ as $t\downarrow0$,
and the potentials on the left converge uniformly to $v$.
Integrate against nonnegative continuous tests and pass to the limit.
Subtract the first restricted measure to obtain the claimed contact
inequality. Replacing one factor at a time proves the pure-power claim
and the analogous contact inequalities for all mixed powers.
\end{proof}

For a continuous real obstacle $h$, define $P(h)$ to be the supremum of
continuous admissible functions bounded above by $h$. Constants make the
family nonempty. The penalty construction also proves that this
supremum is continuous and its Hessian measure lies on the contact set.
\index{Envelope!continuous Hessian obstacle}

\begin{lemma}\label{hessian:continuous-obstacle-envelope}
The function $P(h)$ is continuous and admissible, and $H_m(P(h))$ is
supported on $\{P(h)=h\}$. Moreover
$\|P(h_1)-P(h_2)\|_\infty\leq\|h_1-h_2\|_\infty$.
\end{lemma}
\begin{proof}
For smooth $h$, use the penalty solutions $w_N$ in \cref{hessian:continuous-smooth-approximation}. Its maximum argument applies to every continuous admissible
competitor $v\leq h$, with an error depending on $h$ and $N$, and not
on $v$. Taking their supremum gives
\[
 P(h)-E_N-2M_N/N\leq w_N\leq P(h),
\]
where the error tends to zero. Thus $P(h)$ is their uniform limit and
is continuous and admissible. On every compact subset of
$\{P(h)<h\}$ the penalty densities
$b e^{N(w_N-h)}$ tend uniformly to zero. Product continuity implies
the contact support assertion.
Adding constants to competitors gives the stated Lipschitz bound.
Approximate a continuous $h$ uniformly by smooth obstacles. Their
envelopes converge uniformly to $P(h)$, proving continuity and
admissibility. Any compact subset of $\{P(h)<h\}$ also lies outside the
contact sets of all sufficiently close smooth obstacles; product
continuity gives the support assertion for $h$.
\end{proof}

The obstacle also bounds discontinuous competitors. Their upper-test
admissibility suffices for the penalty estimate, so allowing them does
not change the continuous envelope of a smooth obstacle.

\begin{proposition}\label{hessian:bounded-envelope}
For every smooth real obstacle $h$, the envelope $P(h)$ equals the
pointwise supremum of all bounded admissible functions below $h$.
\end{proposition}
\begin{proof}
The solutions $w_N$ of \cref{hessian:bounded-penalty-squeeze} are
continuous admissible competitors, so $w_N\leq P(h)$. For each
bounded admissible $u\leq h$, the same lemma gives
\[
 u\leq w_N+E_N+2M_N/N\leq P(h)+E_N+2M_N/N.
\]
Let $N\to\infty$ for this fixed obstacle and fixed $b$.
Thus $P(h)$ dominates every bounded competitor. Since $P(h)$ itself
is a bounded competitor, it equals their supremum.
\end{proof}

The trace test in \eqref{eq:hessian-continuous-admissibility} gives
$L_\omega u\geq-n$, where
$L_\omega u=n\ddc u\wedge\omega^{n-1}/\omega^n$.
We use the classical normalized Green estimate
\begin{equation}\label{eq:hessian-normalized-l1}
 \int_X(\sup_Xu-u)\omega^n\leq C_\omega
\end{equation}
for continuous admissible $u$. The stationary probability of $L_\omega$
has a smooth positive density equivalent to $\omega^n$. The positive
measure $L_\omega u+n$ has stationary mass $n$, so the Green-kernel
estimate gives this bound. It is a linear elliptic input, and does not
assert compactness for a nonlinear Hessian measure.

To use envelopes in domination, we must first exclude a potential with
zero total Hessian mass.

\begin{lemma}\label{hessian:hessian-nonvanishing}
Every continuous admissible $u$ satisfies $\int_X H_m(u)>0$.
\end{lemma}
\begin{proof}
Suppose this integral is zero. For $s>0$ put
$w=\max(u,\min_Xu+s)$ and $U=\{u<\min_Xu+s\}$.
The open set $U$ is nonempty, $0\leq w-u\leq s$, and $w=u$ outside it.
Set $\sigma_u=\omega+\ddc u$, $\sigma_w=\omega+\ddc w$, and
$\mu_q=\sigma_u^q\wedge\sigma_w^{m-q}\wedge\omega^{n-m}$.
The contact lemma gives $\mathbf1_{X\setminus U}\mu_q\geq
\mathbf1_{X\setminus U}\mu_m$, for $0\leq q\leq m$.
For $q<m$ the telescoping identity is
\[
 \mu_q-\mu_m=\ddc(w-u)\wedge T_q,\qquad
 T_q=\sigma_u^q\wedge
       \sum_{r=0}^{m-q-1}\sigma_w^r\wedge\sigma_u^{m-q-1-r}
                                      \wedge\omega^{n-m}.
\]
Integration by parts on $X$ therefore gives
\[
 \int_U(\mu_q-\mu_m)\leq\int_X(\mu_q-\mu_m)
                                =\int_X(w-u)\ddc T_q.
\]
This identity follows first for smooth approximants and then by uniform
product continuity. The derivatives of $T_q$ contain only fixed metric
derivatives and products of at most $m-1$ admissible factors, so the
absolute-form lemma makes $\ddc T_q$ a measure of uniformly controlled
variation. Uniform convergence also justifies its continuous weight.
On $U$, locality gives $\sigma_w=\omega$ and thus
\[
 |\ddc T_q|\leq C\sum_{r=0}^{m-1}\mu_r.
\]
The constant depends only on the metric and $m,n$, and not on $u$ or $s$.
Consequently
\[
 \sum_{q=0}^m\int_U\mu_q
 \leq(m+1)\int_U H_m(u)+Cms\sum_{r=0}^{m-1}\int_U\mu_r.
\]
For sufficiently small $s$, absorption and $H_m(u)=0$ force all these
nonnegative integrals to vanish. But $\mu_0=\omega^n$ on the nonempty
open set $U$, a contradiction.
\end{proof}

The envelope of a stretched difference now converts vanishing of a
restricted measure into an order relation. The linear bound
\eqref{eq:hessian-normalized-l1} prevents these envelopes from tending to
$-\infty$ on an open set.

\begin{proposition}\label{hessian:zero-hessian-domination}
If continuous admissible $u,v$ satisfy
$\mathbf1_{\{u<v\}}H_m(u)=0$, then $u\geq v$.
\end{proposition}
\begin{proof}
Replace $v$ by $\max(u,v)$, so $u\leq v$.
For $b>1$ put $p_b=P(bu-(b-1)v)$ and
$z_b=b^{-1}p_b+(1-b^{-1})v\leq u$.
On the contact set $C_b=\{p_b=bu-(b-1)v\}$ one has $z_b=u$;
the contact inequality and polarized positivity give
\[
 \mathbf1_{C_b}b^{-m}H_m(p_b)
                 \leq\mathbf1_{C_b}H_m(z_b)
                 \leq\mathbf1_{C_b}H_m(u).
\]
The envelope measure is supported on $C_b$ and has positive mass.
It vanishes on $C_b\cap\{u<v\}$ by the hypothesis, so some point of
$C_b\cap\{u=v\}$ exists. Hence
$\sup_Xp_b\geq\inf_Xv$, while $p_b\leq u$.
The estimate \eqref{eq:hessian-normalized-l1} is consequently
uniform for $p_b$. On the open set $E_a=\{u<v-a\}$, for $a>0$, one has
$p_b\leq v-ba$, and therefore
\[
 C_\omega\geq\int_X(\sup_Xp_b-p_b)\omega^n
                     \geq(ba-\operatorname{osc}_Xv)\int_{E_a}\omega^n.
\]
Letting $b\to\infty$ gives zero volume for $E_a$. Since it is open it
is empty, and $a\downarrow0$ gives $u\geq v$.
\end{proof}

\begin{proposition}\label{hessian:strict-hessian-domination}
If continuous admissible $u,v$ satisfy
\[
 \mathbf1_{\{u<v\}}H_m(u)\leq
                    c\mathbf1_{\{u<v\}}H_m(v),\qquad 0\leq c<1,
\]
then $u\geq v$.
\end{proposition}
\begin{proof}
Choose $b>1$ so that $(1-b^{-1})^m>c$, and put
$p=P(bu-(b-1)v)$ and $z=b^{-1}p+(1-b^{-1})v\leq u$.
On $C=\{p=bu-(b-1)v\}$ the contact comparison gives
\[
 \mathbf1_C\bigl(b^{-m}H_m(p)+(1-b^{-1})^mH_m(v)\bigr)
                                  \leq\mathbf1_C H_m(u).
\]
Restricting to $\{u<v\}$ and using the hypothesis shows that
$H_m(p)$ vanishes on $C\cap\{u<v\}$. Its support lies in $C$, where
$p<v$ is equivalent to $u<v$. Thus $H_m(p)$ vanishes on $\{p<v\}$.
The preceding domination proposition gives $p\geq v$.
Since $p\leq bu-(b-1)v$, it follows that $u\geq v$.
\end{proof}

Strict domination yields the comparison needed for the proper equations.
\index{Domination principle!continuous Hessian measures}

\begin{lemma}\label{hessian:hessian-exponential-order}
For continuous admissible $u,v$ and $a>0$, if
\[
 e^{-av}H_m(v)\geq e^{-au}H_m(u),
\]
then $u\geq v$.
\end{lemma}
\begin{proof}
Fix $\delta>0$. On $\{u<v-\delta\}$, the assumed inequality gives
\[
 H_m(u)\leq e^{a(u-v)}H_m(v)
       \leq e^{-a\delta}H_m(v-\delta).
\]
The constant $e^{-a\delta}$ is strictly smaller than one, so \cref{hessian:strict-hessian-domination} applied to $u$ and $v-\delta$ gives $u\geq v-\delta$. Letting $\delta\downarrow0$ proves the assertion.
\end{proof}

\section{Regularity of continuous solutions}
\label{sec:hessian-continuous-regularity}
\index{Hessian equation!regularity of continuous solutions}

The proper equation is stable in the uniform norm, so its smooth solutions
extend to continuous data. This also identifies the limit of the
auxiliary equations with a prescribed continuous weak solution.

\begin{proposition}\label{hessian:proper-uniform-stability}
Let $H_j$ be smooth and converge uniformly to a continuous $H$.
The smooth solutions $w_j$ of
$H_m(w_j)=e^{w_j+H_j}\omega^n$ converge uniformly to the unique
continuous admissible solution of
$H_m(w)=e^{w+H}\omega^n$. For smooth $H_1,H_2$,
\[
 \|w_1-w_2\|_\infty\leq\|H_1-H_2\|_\infty.
\]
\end{proposition}
\begin{proof}
At a maximum of $w_1-w_2$, monotonicity gives
$e^{w_1+H_1}\leq e^{w_2+H_2}$, hence the maximum is at most
$\|H_1-H_2\|_\infty$. Exchange the functions to obtain the norm bound.
Thus $w_j$ is Cauchy uniformly. The limit is admissible by passage of
the linear distributional inequalities to the limit. \Cref{hessian:continuous-products} passes the equation to this limit. \Cref{hessian:hessian-exponential-order} proves uniqueness among continuous
admissible solutions, independently of the approximation of $H$.
\end{proof}

For smooth positive densities depending on the potential, the estimate
of \cref{sec:hessian-trace} gives complex-Hessian bounds independent of
derivatives of the prescribed approximation. The elliptic extension then
upgrades the continuous solution to a classical one.

\Needspace{11\baselineskip}
\begin{theorem}\label{hessian:continuous-hessian-regularity}
If $\varphi$ is continuous and admissible on $X$, $f(x,t)$ is smooth
on a neighborhood of its bounded range, and
\[
 H_m(\varphi)=e^{f(x,\varphi)}\omega^n,
\]
then $\varphi$ is smooth.
\end{theorem}
\begin{proof}
Choose smooth admissible $v_j\to\varphi$ uniformly, by \cref{hessian:continuous-smooth-approximation}. Their values lie in a fixed compact interval on which $f$ is
smooth. Put $h(x,t)=-t+f(x,t)$ and solve
\[
 H_m(u_j)=e^{u_j+h(x,v_j)}\omega^n.
\]
The preceding stability proposition shows that $u_j$ converges uniformly
to a continuous admissible $u$ satisfying
\[
 e^{-u}H_m(u)=e^{-\varphi+f(x,\varphi)}\omega^n
                              =e^{-\varphi}H_m(\varphi).
\]
\Cref{hessian:hessian-exponential-order}, applied in both directions, gives $u=\varphi$.

The zero-order bounds for $u_j$ are uniform. Fix $M$ bounding them
and choose $B>\|\varphi\|_\infty+1$. Extend $h$ smoothly away from a neighborhood of the range of $\varphi$ to $X\times[-B-1,B+1]$; this does not change the equations for large $j$. The tolerance $\varepsilon$ in
\cref{hessian:prescribed-potential-trace} depends on these bounds and on
$h$, without derivatives of the smooth comparison function.
Choose a fixed smooth $\widetilde\varphi$ with
$\|\widetilde\varphi-\varphi\|_\infty<\varepsilon/2$ and
$\|\widetilde\varphi\|_\infty\leq B$.
For large $j$, $\|v_j-\widetilde\varphi\|_\infty<\varepsilon$.
\Cref{hessian:prescribed-potential-complex-hessian} then gives uniform gradient
and complex-Hessian bounds for $u_j$, independent of derivatives of
$v_j$. \Cref{hessian:complex-hessian-compactness} yields
$\varphi\in W^{2,p}$ for every finite $p$, and
$\varphi\in C^{1,\alpha}$ for every $\alpha<1$.

\Cref{hessian:proper-auxiliary-limit} makes $\varphi$ a continuous viscosity
solution of the concave uniformly elliptic extension, with right-hand
side $e^{f(x,\varphi)/m}$. This right-hand side is H\"older. The analytic estimate in \cref{sec:hessian-complex-limits} therefore gives
$\varphi\in C^{2,\beta}$, for some $\beta>0$.
The operator agrees with the original Hessian root on the attained
compact cone set; the now classical equation thus has smooth dependence
on $x,\varphi$ and its complex Hessian, and is uniformly elliptic there.
Difference quotients and interior Schauder estimates give $C^{3,\beta}$;
iteration proves smoothness. The estimate applies directly to the viscosity limit.
\end{proof}

Continuity of the given solution is a hypothesis of this theorem.
The passage from an arbitrary bounded admissible weak solution to this
continuous regime requires a product theory and a continuity estimate
for that larger class.

\section{Bounded potentials and recursive Hessian currents}
\label{sec:hessian-bounded-currents}\index{Hessian measure!bounded recursive currents}

The decreasing approximation of bounded admissible potentials supplies
limits that uniform approximation cannot reach. A bounded potential can
multiply a current with measure coefficients, so one can apply $\ddc$
to the resulting current. Repeating this operation requires bounds on
every coefficient, rather than only on the trace of a positive mixed
measure. The absolute-form estimate provides these bounds below degree
$m$; at degree $m$, positivity is obtained after wedging with
$\omega^{n-m}$.

All bounded potentials have their canonical upper semicontinuous representatives from
\eqref{eq:hessian-bounded-upper-tests}. Their pointwise values will be
used when multiplying measure coefficients. For a real current $T$ of
order zero, write $\|T\|_K$ for the sum of the total variations of its
coefficients on a compact coordinate set $K$. Fix $K\Subset U$ in a
coordinate chart. Different fixed coordinate frames give equivalent
local bounds.

\Needspace{11\baselineskip}
\begin{lemma}\label{hessian:bounded-zero-order}
If $u$ is bounded admissible, then $\ddc u$ has signed Radon measure
coefficients. For $K\Subset U$ contained in a coordinate chart,
\[
 \|\ddc u\|_K\leq C_{K,U,\omega}(1+\|u\|_{L^\infty(X)}).
\]
\end{lemma}
\begin{proof}
Choose smooth strictly admissible $u_j\downarrow u$ by
\cref{hessian:bounded-decreasing-approximation}, and put
$\sigma_j=\omega+\ddc u_j$. After discarding finitely many terms,
$\|u_j\|_\infty\leq\|u\|_\infty+1$. Indeed, the increasing open sets
$\{u_j<\sup_Xu+1\}$ cover $X$, and compactness gives a single index
whose open set is all of $X$; the lower bound follows from $u_j\geq u$.

The closed $m$-cone lies in the closed $2$-cone. In an $\omega$-unitary
frame, its eigenvalues satisfy
\[
 \sum_i\lambda_i^2\leq(\operatorname{tr}\sigma_j)^2,
 \qquad \operatorname{tr}\sigma_j\geq0.
\]
Thus every coefficient of $\sigma_j$ is bounded by its trace. The
degree-one Chern--Levine--Nirenberg estimate in
\cref{hessian:continuous-products}, applied to these smooth functions,
bounds the integral of that trace on $K$ by
$C(1+\|u\|_\infty)$. Subtracting the fixed smooth $\omega$ gives the
same bound for the coefficient variations of $\ddc u_j$.
Bounded convergence gives $u_j\to u$ in local $L^1$, and therefore
$\ddc u_j\to\ddc u$ in distributions. The uniform coefficient bounds
show that this limit has measure coefficients with the asserted bound.
\end{proof}

Set $D_0=1$. For an ordered list $u_1,\ldots,u_p$ of bounded
admissible potentials, define recursively
\begin{equation}\label{eq:hessian-bounded-recursion}
 D_p(u_1,\ldots,u_p)
 =\ddc\bigl(u_pD_{p-1}(u_1,\ldots,u_{p-1})\bigr).
\end{equation}
The next result justifies the bounded multiplication at every step up
to $p=m$.

\Needspace{13\baselineskip}
\begin{proposition}\label{hessian:bounded-recursive-currents}
For $1\leq p<m$, the recursion defines a real closed current of order
zero and bidegree $(p,p)$. If
$M=\max_a\|u_a\|_{L^\infty(X)}$, then
\[
 \|D_p(u_1,\ldots,u_p)\|_K
       \leq C_{K,U,\omega,p}(1+M)^p.
\]
For $p=m$, it defines a real closed distributional current. For any
smooth admissible sequences $u_a^{j_a}\downarrow u_a$ with a common
uniform bound, one has
\begin{equation}\label{eq:hessian-bounded-iterated-limit}
 D_p(u_1,\ldots,u_p)
 =\lim_{j_p\to\infty}\cdots\lim_{j_1\to\infty}
    \ddc u_1^{j_1}\wedge\cdots\wedge\ddc u_p^{j_p}
\end{equation}
in distributions. The index $j_1$ tends to infinity first. Below degree
$m$, all these limits have locally uniformly bounded coefficient
variations.
\end{proposition}
\begin{proof}
The case $p=1$ is the preceding lemma. Suppose the assertions hold
through degree $p-1<m$. Since $u_p$ is bounded Borel and $D_{p-1}$
has measure coefficients, their product is well defined. Applying
$\ddc$ gives a real current that is both $\mathrm d$- and
$\dc$-closed.

For a fixed smooth complementary test form $\eta$, dominated
convergence against the fixed signed measures gives
\[
 \begin{aligned}
 \langle D_p,\eta\rangle
 &=\langle u_pD_{p-1},\ddc\eta\rangle\\
 &=\lim_{j_p\to\infty}
       \langle u_p^{j_p}D_{p-1},\ddc\eta\rangle.
 \end{aligned}
\]
With $j_p$ fixed, the inductive iterated limit for $D_{p-1}$ can be
paired with the smooth form $u_p^{j_p}\ddc\eta$. Integration by parts
for the smooth closed factors identifies the resulting pairing with
the smooth product on the right of
\eqref{eq:hessian-bounded-iterated-limit}. This proves that identity
also when $p=m$.

For $p<m$, expand the smooth product using
$\ddc u_a^{j_a}=\sigma_a^{j_a}-\omega$. A term with $q$ admissible
factors, $0\leq q\leq p$, has the form
\[
 \sigma_{a_1}^{j_{a_1}}\wedge\cdots\wedge
 \sigma_{a_q}^{j_{a_q}}\wedge\omega^{p-q}.
\]
Every such $q$ is strictly less than $m$. Apply
\cref{hessian:absolute-mixed-forms} to each coefficient test, and then
the smooth mass estimate from \cref{hessian:continuous-products} to
the potential sum. The term $q=0$ is the fixed smooth form $\omega^p$.
For $q\geq1$, the sum of its admissible forms is
$q(\omega+\ddc\bar u)$, where $\bar u$ is their average and
$\|\bar u\|_\infty\leq M+1$. The result is a
coefficient variation bound $C(1+M)^p$, after choosing the
approximations with norm at most $M+1$ as in the preceding proof.
For arbitrary uniformly bounded approximations the same argument
uses their common bound. The bounds survive each distributional
limit, proving that $D_p$ has order zero and the stated estimate.
This completes the induction. No coefficient bound for the unwedged
degree-$m$ current is used or asserted.
\end{proof}

\Needspace{8\baselineskip}
For a subset $I=\{i_1<\cdots<i_q\}\subset\{1,\ldots,p\}$, write
$D_I=D_q(u_{i_1},\ldots,u_{i_q})$, with $D_\varnothing=1$.
Define the shifted current and its mixed positive current by
\begin{equation}\label{eq:hessian-bounded-shifted-current}
 R_p(u_1,\ldots,u_p)
 =\sum_{I\subset\{1,\ldots,p\}}\omega^{p-|I|}\wedge D_I,
 \qquad Q_p=R_p\wedge\omega^{n-m}.
\end{equation}
Below degree $m$, $R_p$ has measure coefficients. At degree $m$ it
is initially a distributional current. Its positive wedge $Q_m$ is
nevertheless a measure.

Put $M=\max_a\|u_a\|_{L^\infty(X)}$.

\Needspace{13\baselineskip}
\begin{theorem}\label{hessian:bounded-shifted-positive}\index{Hessian measure!iterated decreasing limits}
For $1\leq p\leq m$, $Q_p$ is a positive current of order zero, and
\[
 \int_KQ_p\wedge\omega^{m-p}
       \leq C_{K,U,\omega,p}(1+M)^p.
\]
It is the iterated limit, in the order of
\eqref{eq:hessian-bounded-iterated-limit}, of
\[
 (\omega+\ddc u_1^{j_1})\wedge\cdots\wedge
 (\omega+\ddc u_p^{j_p})\wedge\omega^{n-m}.
\]
In particular $Q_m$ is a positive Radon measure. These currents are
local in the canonical potentials and independent of the chosen
approximations for the fixed ordering. When all potentials are
continuous, $Q_p$ equals the continuous mixed product of
\cref{hessian:continuous-products}.
\end{theorem}
\begin{proof}
Expand the smooth shifted product over its subsets. Applying
\eqref{eq:hessian-bounded-iterated-limit} to each subset, in its
inherited ordering, gives \eqref{eq:hessian-bounded-shifted-current}.
The absent indices have no effect on a subset's limit.

Each smooth product is positive: pairing with a decomposable
strongly positive complementary test appends $m-p$ nonnegative
rank-one $(1,1)$-forms, and polarized G\aa rding positivity applies.
The smooth mixed mass bound is uniform in all indices. Distributional
limits preserve positivity and the local trace bound at every step.
Positive currents have order zero, so this gives the assertions for
$Q_p$, including the top-degree measure.

The recursion uses only bounded multiplication of canonical values
and local distributional differentiation. It is therefore local.
Its definition contains no approximation choices, so the proved
limit is independent of those choices while the ordering is fixed.
For continuous potentials, the approximations converge uniformly by
Dini's theorem. The telescoping estimate in
\cref{hessian:continuous-products} is uniform in all remaining
indices and shows that their iterated limit is the same continuous
mixed product. This also proves permutation invariance in the
continuous case.
\end{proof}

The shifted currents are not closed. Their derivatives can be stated
without multiplying a discontinuous function by a distribution of
positive order. This is the torsion identity needed before considering
integration by parts with bounded weights.

\Needspace{15\baselineskip}
\begin{proposition}\label{hessian:bounded-current-torsion}\index{Torsion!bounded Hessian currents}
Write $R_{p-1}^{\widehat a}$ for the shifted current obtained by
omitting $u_a$, in the inherited order, and use two hats for two
omissions. Put $R_0=1$. For $1\leq p\leq m$,
\[
 \begin{aligned}
 \mathrm dR_p&=\mathrm d\omega\wedge
                         \sum_{a=1}^pR_{p-1}^{\widehat a},\\
 \dc R_p&=\dc\omega\wedge
                         \sum_{a=1}^pR_{p-1}^{\widehat a},\\
 \ddc R_p&=\ddc\omega\wedge
                         \sum_{a=1}^pR_{p-1}^{\widehat a}
       +2\,\mathrm d\omega\wedge\dc\omega\wedge
                         \sum_{a<b}R_{p-2}^{\widehat a\widehat b}.
 \end{aligned}
\]
The last sum is absent for $p=1$. All currents on the right have order
zero. If all but at most one potential are smooth, $R_p$ is independent
of the ordering and is the ordinary
product of the smooth factors with $\omega+\ddc u$.
\end{proposition}
\begin{proof}
Every $D_I$ is both $\mathrm d$- and $\dc$-closed. Differentiating
\eqref{eq:hessian-bounded-shifted-current} therefore gives
\[
 \mathrm dR_p=\mathrm d\omega\wedge
 \sum_I(p-|I|)\omega^{p-|I|-1}\wedge D_I,
\]
with zero terms omitted. The coefficient $p-|I|$ counts the omitted
indices outside $I$, which proves the first formula; the second is
identical with $\dc$ in place of $\mathrm d$. Apply $\mathrm d$ to
the second formula. Each pair of omitted indices occurs twice, and
$-\dc\omega\wedge\mathrm d\omega
=\mathrm d\omega\wedge\dc\omega$ because both forms have odd
degree three. This gives the displayed $\ddc$ formula, with every
background derivative retained.

If all but one potential are smooth, applying $\ddc$ to the bounded
potential times a smooth closed product gives its $\ddc$ wedged
with that product. At later smooth steps, closedness again makes
the recursion ordinary smooth multiplication. This is independent
of the position of the bounded potential, and the subset expansion
gives the last assertion.
\end{proof}

The iterated limit fixes the order of the potentials. To compare it
with a simultaneous limit, discontinuous weights must be controlled
against varying coefficient measures. The next section supplies this
control for the first Hessian, giving symmetry and simultaneous
decreasing convergence for two factors in
\cref{hessian:bounded-two-factor-convergence}. Higher-degree products
still need additional estimates.

\section{Sobolev capacity and two Hessian factors}
\label{sec:hessian-sobolev-two-factors}

Weak convergence of measures is not enough to pass to the limit in
their products with bounded potentials; one needs a capacity that
dominates the measures and makes the potentials quasicontinuous.
Sobolev capacity does this for the first Hessian of a bounded
admissible potential, through its linear trace, but a product of two
Hessians can charge a set of zero Sobolev capacity. The trace is taken
against a $\ddc$-closed conformal multiple of $\omega^{n-1}$: an
admissible form need not be positive in every tangent direction, so
only its trace relative to $\omega$ is known to be nonnegative.

Choose $\rho>0$ by \cref{thm:gauduchon} for $\omega$ and put
\[
 T=\rho\omega^{n-1},\qquad \ddc T=0,\qquad
 A=\int_X\omega\wedge T>0.
\]
Write $H^1(X)=W^{1,2}(X,\mathbb R)$ and use the equivalent Sobolev norm
\[
 \|f\|_{H^1}^2=\int_Xf^2\omega^n+
                   \int_X\mathrm df\wedge\dc f\wedge T.
\]
For an open set $G\subset X$ define
\[
 \operatorname{Cap}_1(G)
 =\inf\{\|f\|_{H^1}^2:f\in H^1(X),\ f\geq1
                          \text{ almost everywhere on }G\},
\]
and for an arbitrary set $E$ take the infimum over open $G\supset E$.
It is distinct from the local Hessian capacities of
\cref{sec:hessian-smooth-capacity}. Truncation allows $0\leq f\leq1$ in the infimum for open sets.
\index{Capacity!Sobolev}

\begin{lemma}\label{hessian:sobolev-capacity-tools}
The capacity $\operatorname{Cap}_1$ is countably subadditive. For a
smooth real $f$ and $a>0$,
\[
 \operatorname{Cap}_1(\{|f|>a\})\leq a^{-2}\|f\|_{H^1}^2.
\]
If a positive Radon measure $\mu$ satisfies
$|\int_X\phi\,\mathrm d\mu|\leq C\|\phi\|_{H^1}$ for every smooth real
$\phi$, then for every Borel set $E$,
\[
 \mu(E)\leq C\operatorname{Cap}_1(E)^{1/2}.
\]
\end{lemma}
\begin{proof}
Let $G_i$ be open sets with $\sum_i\operatorname{Cap}_1(G_i)<\infty$;
if the sum is infinite there is nothing to prove. Fix $\epsilon>0$ and
choose Sobolev functions $0\leq f_i\leq1$ equal to one almost
everywhere on $G_i$, with
$\|f_i\|_{H^1}^2\leq\operatorname{Cap}_1(G_i)+\epsilon2^{-i}$.
The Sobolev chain rule gives
\[
 \|\max_{i\leq N}f_i\|_{H^1}^2\leq\sum_{i\leq N}\|f_i\|_{H^1}^2.
\]
Indeed, almost everywhere the gradient of the maximum is the gradient
of an active function, and the squared value is bounded by the sum of
the squared values. These maxima increase with $N$, converge in
$L^2$, are bounded in $H^1$, and have a weak $H^1$ limit equal to one
almost everywhere on $\bigcup_iG_i$. Lower semicontinuity of the norm
gives
$\operatorname{Cap}_1(\bigcup_iG_i)\leq\sum_i\operatorname{Cap}_1(G_i)+\epsilon$.
This proves subadditivity for open sets, and taking open supersets
proves it for arbitrary sets.
The level-set estimate follows by testing with
$\min\{1,|f|/a\}$ and the same Sobolev chain rule.

Fix a compact $K\subset G$ and a Sobolev competitor $0\leq f\leq1$
for the open set $G$. Mollify $f$ in finitely many coordinate charts
by nonnegative kernels and glue the results with a fixed partition of
unity. This gives smooth nonnegative $f_\epsilon\to f$ in $H^1$ that
satisfy $f_\epsilon=1$ on $K$ for small $\epsilon$: every convolution
ball centered on $K$ then lies in $G$, where $f=1$ almost everywhere.
The assumed estimate gives
$\mu(K)\leq C\lim_\epsilon\|f_\epsilon\|_{H^1}=C\|f\|_{H^1}$.
Take the infimum over $f$ and then the supremum over compact $K\subset G$
to obtain the estimate for $G$. For a Borel set $E$, apply it to
every open $G\supset E$ and use $\mu(E)\leq\mu(G)$.
\end{proof}

\begin{lemma}\label{hessian:bounded-sobolev-trace}\index{Sobolev regularity!bounded Hessian potentials}
Every bounded admissible $u$ belongs to $H^1(X)$. Its trace
\[
 \mu_u=(\omega+\ddc u)\wedge T
\]
is a positive Radon measure of total mass $A$. If $\|u\|_\infty\leq M$,
then
\[
 \|u\|_{H^1}\leq C(1+M),\qquad
 \mu_u(E)\leq C(1+M)\operatorname{Cap}_1(E)^{1/2}
\]
for every Borel $E$. On each fixed relatively compact coordinate chart,
the total variation of every coefficient of $\ddc u$ is also bounded
on $E$ by $C(1+M)\operatorname{Cap}_1(E)^{1/2}$.
\end{lemma}
\begin{proof}
First let $u$ be smooth admissible. The first-cone inequality gives
$\mu_u\geq0$, since $T$ is conformal to $\omega^{n-1}$, and
$\ddc T=0$ gives $\mu_u(X)=A$. Stokes' theorem and $\ddc T=0$ give
$\int_X\ddc(u^2)\wedge T=0$. Since
$\ddc(u^2)=2u\ddc u+2\mathrm du\wedge\dc u$, this yields the energy
identity
\[
 \int_X\mathrm du\wedge\dc u\wedge T
   =-\int_Xu\,\ddc u\wedge T
   =-\int_Xu\,\mathrm d\mu_u+\int_Xu\omega\wedge T
   \leq2MA.
\]
Together with the bound $M^2\int_X\omega^n$ for the $L^2$ term, this
gives $\|u\|_{H^1}\leq C(1+M)$.
For bounded admissible $u$, choose smooth admissible $u_j\downarrow u$
by \cref{hessian:bounded-decreasing-approximation}, discard a finite
initial segment, and take $\|u_j\|_\infty\leq M+1$ as in
\cref{hessian:bounded-zero-order}. Bounded convergence gives strong
$L^2$ convergence. The energy bound gives weak $H^1$ compactness, and
every weak cluster point agrees with $u$ in $L^2$; thus the entire
sequence converges weakly to $u$ in $H^1$. Lower semicontinuity of the
norm proves the Sobolev bound.
The trace measures converge weakly to $\mu_u$, preserving positivity and total mass.

For any smooth real $\phi$, the first-order Stokes identity is
\[
 \int_X\phi\,\ddc u\wedge T
  =-\int_X\mathrm d\phi\wedge\dc u\wedge T
       +\int_X\phi\,\dc u\wedge\mathrm dT.
\]
It remains valid for $u\in H^1$ by smooth Sobolev approximation.
Both terms are bounded by $C\|u\|_{H^1}\|\phi\|_{H^1}$; the background term
$\int\phi\omega\wedge T$ is bounded by $C\|\phi\|_{H^1}$.
Hence $|\int_X\phi\,\mathrm d\mu_u|\leq C(1+M)\|\phi\|_{H^1}$, and
\cref{hessian:sobolev-capacity-tools} gives the bound for $\mu_u$.

For a smooth admissible approximant, the eigenvalues of
$\omega+\ddc u_j$ relative to $\omega$ have modulus at most its trace,
by the proof of \cref{hessian:weighted-cone-trace}, since $m\geq2$.
Consequently every coefficient of this form is dominated by a fixed
multiple of its trace.
On a fixed coordinate chart this bounds the coefficient variations of
$\ddc u_j$ by $C(\mu_{u_j}+\omega^n)$. Passing to the weak measure
limits of \cref{hessian:bounded-zero-order} gives the same domination
by $C(\mu_u+\omega^n)$. Finally
$\int_E\omega^n\leq\operatorname{Cap}_1(E)$ and
$\operatorname{Cap}_1(E)\leq\|1\|_{H^1}^2$ absorb the smooth term into the
claimed square-root estimate.
\end{proof}

\begin{lemma}\label{hessian:bounded-sobolev-quasicontinuity}\index{Admissible potential!Sobolev quasicontinuity}
Every bounded admissible potential is continuous on the
complement of an open set of arbitrarily small $\operatorname{Cap}_1$.
The same open set can be chosen for any countable family of such
potentials.
\end{lemma}
\begin{proof}
Choose uniformly bounded smooth admissible $u_j\downarrow u$ by
\cref{hessian:bounded-decreasing-approximation}. By the preceding lemma
these functions converge weakly to $u$ in $H^1$. Mazur's lemma gives
finite convex combinations $w_N$ of terms with $j\geq N$ such that
\[
 \|w_N-u\|_{H^1}\leq8^{-N}.
\]
Because $u\leq w_N\leq u_N$, they converge pointwise to the canonical
function $u$, not just to an unspecified almost-everywhere representative.
Put
\[
 G_{N_0}=\bigcup_{N\geq N_0}
                  \{|w_{N+1}-w_N|>2^{-N}\}.
\]
\Cref{hessian:sobolev-capacity-tools} implies
\[
 \operatorname{Cap}_1(G_{N_0})
 \leq\sum_{N\geq N_0}4^N\|w_{N+1}-w_N\|_{H^1}^2
 \leq C\sum_{N\geq N_0}16^{-N}\longrightarrow0.
\]
Outside $G_{N_0}$, the successive differences form a uniformly summable
series. Thus $w_N$ converge uniformly there and their limit $u$ is
continuous on that closed set. For a countable family, choose the
exceptional open sets with capacities less than $\epsilon2^{-i}$ and
take their union. Countable subadditivity proves the last assertion.
\end{proof}

The capacity estimate depends on using a first Hessian. A product
of two Hessians can already charge a set of zero Sobolev capacity.

\begin{lemma}\label{hessian:sobolev-codimension-two}\index{Capacity!Sobolev zero sets}
Let $E$ be a compact subset of a smooth real submanifold of $X$ of
codimension two. Then $\operatorname{Cap}_1(E)=0$.
\end{lemma}
\begin{proof}
It suffices to treat a compact subset of a submanifold chart, since
finitely many such subsets cover $E$ and the capacity is subadditive.
Choose real coordinates $(s,t)$ there, with $t\in\mathbb R^2$ and the
submanifold given by $t=0$. Choose $0\leq\eta\leq1$ smooth,
supported in this chart and equal to one near the compact subset.
Put $r=|t|$.
For $0<\epsilon<1/2$ define
\[
 h_\epsilon(r)=
 \begin{cases}
  1,&0\leq r\leq\epsilon^2,\\
  \displaystyle\frac{\log(\epsilon/r)}{\log(1/\epsilon)},
       &\epsilon^2<r<\epsilon,\\
  0,&r\geq\epsilon.
 \end{cases}
\]
The Sobolev function $f_\epsilon=\eta h_\epsilon$ equals one on an
open neighborhood of that subset. On the fixed compact coordinate
set the Sobolev norm is equivalent to the Euclidean one. Integration
in the two normal variables gives
\[
 \|f_\epsilon\|_{H^1}^2
 \leq C\left(\epsilon^2+
     \int_{\epsilon^2}^{\epsilon}
       \frac{\mathrm dr}{r\log(1/\epsilon)^2}\right)
 =C\left(\epsilon^2+\frac1{\log(1/\epsilon)}\right)
 \longrightarrow0.
\]
The derivatives of $\eta$ contribute at most $C\epsilon^2$.
Since $f_\epsilon$ is a competitor for an open neighborhood of the
subset, the capacity of the subset is at most
$\|f_\epsilon\|_{H^1}^2$ for every $\epsilon$, which proves the
assertion.
\end{proof}

\begin{example}\label{hessian:sobolev-two-product-obstruction}\index{Hessian measure!Sobolev capacity obstruction}
There are bounded continuous $\omega$-plurisubharmonic potentials
$u,v$ and a compact set $E$ such that
\[
 \operatorname{Cap}_1(E)=0,\qquad
 \int_E Q_2(u,v)\wedge\omega^{m-2}>0.
\]
Consequently the first-Hessian capacity estimate of
\cref{hessian:bounded-sobolev-trace} does not extend to
this positive two-factor measure, even within the continuous
$\omega$-plurisubharmonic subclass.
\end{example}
\begin{proof}
Fix a holomorphic chart with coordinates $z_j=x_j+\mathrm i y_j$ and a smooth
nonnegative cutoff $\chi$ supported in the chart and equal to one
on a smaller coordinate box. Extend
\[
 u=a\chi|x_1|,\qquad v=a\chi|x_2|
\]
by zero outside the chart. These functions are globally continuous
and bounded. To check their curvature, replace $|x_j|$ by
$g_\delta(x_j)=(x_j^2+\delta^2)^{1/2}$, with $0<\delta\leq1$.
The Hessian of $\chi g_\delta$ has the nonnegative term
$\chi\ddc g_\delta$; all its remaining terms are uniformly bounded
forms, since $g_\delta$ and $g_\delta'$ are uniformly bounded on the
fixed support. Thus $\ddc(\chi g_\delta)\geq-C\omega$ with $C$
independent of $\delta$. Choosing $a>0$ sufficiently small gives
$\omega+\ddc(a\chi g_\delta)\geq\omega/2$ for both coordinates.
Their decreasing limits $u,v$ as $\delta\downarrow0$ are therefore
$\omega$-plurisubharmonic and admissible.

On the box where $\chi=1$, set
$\alpha_j=\frac{\mathrm i}{4\pi}\mathrm dz_j\wedge\mathrm d\bar z_j$.
The book's normalization gives
\[
 \ddc g_\delta(x_j)=\tfrac12g_\delta''(x_j)\alpha_j,
 \qquad g_\delta''(x_j)
       =\frac{\delta^2}{(x_j^2+\delta^2)^{3/2}}
       \longrightarrow2\delta_{\{x_j=0\}}.
\]
The measures $g_\delta^{\prime\prime}(t)\,\mathrm dt$ have mass two,
while their mass outside any fixed neighborhood of zero tends to
zero. Thus $\Theta=\ddc u=a\delta_{\{x_1=0\}}\alpha_1$ on the box. This
current is $\mathrm d$- and $\dc$-closed and depends only on $x_1$,
while $v=a|x_2|$ there. Hence
$\ddc(g_\delta(x_2)\Theta)=\ddc g_\delta(x_2)\wedge\Theta$, and since
$g_\delta\to|\cdot|$ uniformly, \eqref{eq:hessian-bounded-recursion}
gives on the box
\[
 D_2(u,v)=\ddc(v\Theta)
 =\lim_{\delta\downarrow0}a\,\ddc g_\delta(x_2)\wedge\Theta
 =a^2\delta_{\{x_1=0\}}\delta_{\{x_2=0\}}
                         \alpha_1\wedge\alpha_2;
\]
the two factors depend on distinct real variables, so their product
converges to the tensor-product measure.
Choose $E$ to be a compact piece of $\{x_1=x_2=0\}$ inside this
box with nonempty relative interior. The preceding lemma gives
$\operatorname{Cap}_1(E)=0$. On this box, the positive measure
$D_2(u,v)\wedge\omega^{n-2}$ has strictly positive mass on $E$,
since $\omega$ is strictly positive. The smooth term $\omega^n$ and
the two single-Hessian terms in the expansion
\[
 Q_2(u,v)\wedge\omega^{m-2}
 =\omega^n+\omega^{n-1}\wedge(\ddc u+\ddc v)
                         +D_2(u,v)\wedge\omega^{n-2}
\]
have zero mass on $E$. Indeed, on the box $\omega^{n-1}\wedge\ddc u$
is a smooth multiple of the hypersurface measure of $\{x_1=0\}$,
which vanishes on $E$, and similarly for $v$. Hence the asserted
mass is strictly positive. A bound by
$C\operatorname{Cap}_1(F)^b$ for all Borel $F$, with any $b>0$,
would contradict this choice of $E$.
\end{proof}

The preceding Sobolev-capacity estimate therefore cannot be iterated
to control higher-degree coefficient measures. A local Hessian
capacity gives that control in the next section.

\section{Local capacities and decreasing products}
\label{sec:hessian-smooth-capacity}

The codimension-two example
\cref{hessian:sobolev-two-product-obstruction} shows why Sobolev
capacity cannot control a product of Hessians. Local capacities built
from smooth Hessian measures give the stronger control needed for the
recursive currents. Convergence in these capacities, and hence
quasicontinuity, comes from an energy variation that retains the
derivatives of the Hermitian background.

Work on a coordinate ball $U$ with compact closure in a holomorphic
chart, with the restriction of $\omega$. A smooth local
$m$-subharmonic function means a smooth function $f$ on $U$ that is
$m$-subharmonic relative to $\omega$ in the sense defined before
\cref{hessian:local-unshifted-correction}, that is,
$\ddc f\in\overline{\Gamma_m(\omega)}$ pointwise. For such $f$ put
\[
 \Lambda_s(f)=(\ddc f)^s\wedge\omega^{n-s},\qquad0\leq s\leq m,
\]
where $\Lambda_0(f)=\omega^n$. The local mass and gradient estimates
needed for the capacity use ordinary smooth products.

\begin{lemma}\label{hessian:local-smooth-bounds}
Fix $K\Subset U$ and $M<\infty$. For smooth local $m$-subharmonic
$f_1,\ldots,f_q$ with $|f_i|\leq M$, $0\leq q\leq m$, the mass on $K$
of their mixed trace
$\ddc f_1\wedge\cdots\wedge\ddc f_q\wedge\omega^{n-q}$ is at most a
constant depending only on $K,U,\omega,m,M$. If also $f$ is smooth
local $m$-subharmonic with $|f|\leq M$, then, for $q\leq m-1$ and a
constant $C$ depending on the same data,
\[
 \int_K\mathrm df\wedge\dc f\wedge
       \ddc f_1\wedge\cdots\wedge\ddc f_q
                         \wedge\omega^{n-q-1}\leq C.
\]
\end{lemma}
\begin{proof}
For one potential $f$, the form $\omega+\ddc f$ lies in
$\Gamma_m(\omega)$, and $(\omega+\ddc f)^q\wedge\omega^{n-q}$ expands
into nonnegative terms, one of which is $(\ddc f)^q\wedge\omega^{n-q}$.
The Chern--Levine--Nirenberg bound in the proof of
\cref{hessian:continuous-products}, whose proof uses only the values of
the potential on $U$, therefore bounds the mass on $K$ by $C(1+M)^q$.
For mixed factors, positivity of every term in the expansion of
$(\sum_i\ddc f_i)^q\omega^{n-q}$ bounds the mixed trace by a fixed
multiple of the single trace for $\sum_i f_i$. This proves the first
assertion, including its degree-zero convention.

Replace $f$ by $f+M$, so that $0\leq f\leq2M$ without changing its
gradient or Hessian. The identity
\[
 \ddc(f^2)=2f\ddc f+2\mathrm df\wedge\dc f
\]
also shows that $f^2$ is smooth local $m$-subharmonic. Wedge it with
the remaining $q$ Hessians and $\omega^{n-q-1}$. The first term on
the right then becomes nonnegative, since $f\geq0$ and the degree is
$q+1\leq m$. Hence the gradient term is at most the mixed trace of
$f^2,f_1,\ldots,f_q$, whose mass on $K$ is bounded by the first
assertion with $\max\{M,4M^2\}$ in place of $M$.
\end{proof}

\Needspace{16\baselineskip}
\begin{definition}\label{hessian:smooth-local-capacity}\index{Capacity!smooth local Hessian capacity}
For Borel $E\subset U$ define the smooth capacity
\[
 c_s(E)=\sup\left\{\int_E \Lambda_s(a):
        a\in C^\infty(U),\ -1\leq a\leq0,
        \ddc a\in\overline{\Gamma_m(\omega)}\right\},
\]
for $0\leq s\leq m$. Its outer-open version is
\[
 c_s^*(E)=\inf\{c_s(G):E\subset G\subset U,\ G\text{ open}\}.
\]
A function on a compact $K\Subset U$ is quasicontinuous
for this capacity if, for every $\epsilon>0$, it is continuous on
$K\setminus G$ for some open $G\subset U$ with $c_s(G)<\epsilon$.
\end{definition}

In particular $c_0(E)=\int_E\omega^n$. Each $c_s$ is monotone,
countably subadditive and continuous under increasing unions.
For the last assertion, the supremum over competitors commutes with
the increasing supremum over sets. By
\cref{hessian:local-smooth-bounds}, $c_s$ is finite on compact
subsets of $U$. The outer capacity $c_s^*$ is monotone
and countably subadditive and agrees with $c_s$ on open sets.

The capacities control mixed traces for smooth local $m$-subharmonic functions
$f_1,\ldots,f_q$ with a common bound. Choose a fixed smooth strictly
plurisubharmonic $b$ with $\ddc b\geq\omega$ on $U$. For $q\leq s$,
choose constants $B$ and $L\geq1$ so that
$(f_1+\cdots+f_q+b-B)/L$ lies in $[-1,0]$.
Both constants depend only on the common bound and the fixed data.
Positivity of the smooth binomial expansion gives
\begin{equation}\label{eq:hessian:smooth-capacity-bound}
 (\ddc f_1+\cdots+\ddc f_q)^q\omega^{n-q}
       \leq C \Lambda_s\bigl((\textstyle\sum_i f_i+b-B)/L\bigr).
\end{equation}
Indeed, the expansion of the right side contains the term
\[
 (\ddc f_1+\cdots+\ddc f_q)^q(\ddc b)^{s-q}\omega^{n-s}
\]
with a positive coefficient, and replacing any of its $s-q$ factors
$\ddc b$ by $\omega$ subtracts a nonnegative mixed trace.
For $q=0$ the left side is $\omega^n$, and the same argument applies.
Thus every such mixed trace, and every coefficient of a smooth
product of $q<m$ Hessians, is dominated on Borel sets by $C c_s$.
For the coefficients, this follows from the absolute mixed-form
estimate \cref{hessian:absolute-mixed-forms}.

\Needspace{16\baselineskip}
\begin{lemma}\label{hessian:smooth-torsion}\index{Torsion!local Hessian estimates}
Let $0\leq q\leq m-1$ and let $\alpha$ be a smooth real $(1,1)$-form
in $\overline{\Gamma_m(\omega)}$ pointwise.
Put $T=\alpha^q\omega^{n-q-2}$. For smooth real $f$
and nonnegative Borel weights $a$, supported on a fixed compact set,
\[
 \left|\int a\,\mathrm df\wedge\dc\omega\wedge T\right|
 \leq C\left(\int a\,\mathrm df\wedge\dc f\wedge\omega\wedge T\right)^{1/2}
          \left(\int a\,\omega^2\wedge T\right)^{1/2},
\]
where $C$ depends only on $\omega$ and the compact set.
The corresponding inequality with $\dc f\wedge\mathrm d\omega$
holds as well. For a mixed product of $q$ smooth forms in
$\overline{\Gamma_m(\omega)}$ in place of $\alpha^q$, the left side
is bounded by the sum of these bounds over the finitely many
nonempty subset sums of its factors.
\end{lemma}
\begin{proof}
At one point choose an $\omega$-unitary coframe $e_i$
that diagonalizes $\alpha$, with eigenvalues $\lambda\in\overline{\Gamma_m}$. Put $\epsilon_i=\mathrm i e_i\wedge\bar e_i$
and $\Omega=\epsilon_1\cdots\epsilon_n$, so that
$\omega=\sum_i\epsilon_i$ and $\alpha=\sum_i\lambda_i\epsilon_i$.
Writing $B_{ij}=S_{q;ij}(\lambda)$ and $D_i=S_{q;i}(\lambda)$,
ordinary wedge expansion gives
\[
 \alpha^q\omega^{n-q-2}
 =q!(n-q-2)!\sum_{i<j}B_{ij}
                           \prod_{k\notin\{i,j\}}\epsilon_k.
\]
Only the remaining two-plane contributes to its wedge with
$\mathrm df\wedge\dc\omega$. If $\partial f=\sum_i f_i e_i$,
its absolute density relative to $\Omega$ is bounded by
\[
 C\sum_{i<j}(|f_i|+|f_j|)|B_{ij}|.
\]
The constant bounds the tensor coefficients of $\dc\omega$ in all
unitary frames on the fixed compact set; no derivative of this
pointwise frame is taken.

For $q\geq1$, \cref{hessian:deleted-newton} applies, since
$q\leq m-1\leq n-2$, and gives $|B_{ij}|\leq C(D_iD_j)^{1/2}$; for
$q=0$, $B_{ij}=D_i=1$. Finite Cauchy--Schwarz therefore bounds this sum by
\[
 C\sum_i|f_i|\sqrt{D_i}\sum_j\sqrt{D_j}
 \leq C\left(\sum_i|f_i|^2D_i\right)^{1/2}
          \left(\sum_iD_i\right)^{1/2}.
\]
\Needspace{7\baselineskip}
Here $\sum_iD_i=(n-q)S_q(\lambda)$. Moreover
\[
 \begin{aligned}
 \alpha^q\omega^{n-q}&=q!(n-q)!S_q(\lambda)\Omega,\\
 \mathrm df\wedge\dc f\wedge \alpha^q\omega^{n-q-1}
   &=\kappa q!(n-q-1)!\sum_i|f_i|^2D_i\,\Omega,
 \end{aligned}
\]
where $\kappa=1/(2\pi)$ by the book's $\dc$ normalization.
These two densities are nonnegative. Thus the pointwise torsion
density is bounded by a constant times their geometric mean.
Multiplication by $a$ and integral Cauchy--Schwarz give the assertion,
including the cases where either density vanishes. The other wedge
is conjugate up to a fixed sign, which does not affect its absolute
value.

For the mixed assertion with $q\geq1$, use the exact polarization
identity
\[
 q!\,\alpha_1\wedge\cdots\wedge\alpha_q
   =\sum_{J\subset\{1,\ldots,q\}}
           (-1)^{q-|J|}\left(\sum_{i\in J}\alpha_i\right)^q.
\]
Every nonempty subset sum remains in the cone. Applying the single
power estimate to each term and taking absolute values proves the
mixed assertion. At $q=0$ use the single term $\alpha^0=1$.
\end{proof}

\begin{lemma}\label{hessian:boundary-localization}
Let $f_j$ be uniformly bounded smooth local $m$-subharmonic functions
with $f_j\downarrow f$, and fix a smaller concentric ball $V\Subset U$.
There are uniformly bounded smooth local $m$-subharmonic functions
$v_j\downarrow v$ which equal $f_j$ on a neighborhood of $\overline V$
and equal one fixed strictly plurisubharmonic $\psi$ on one fixed
collar $\{R-\delta_0<|z|<R\}$ of the boundary of $U=\{|z|<R\}$.
Moreover $v=f$ near $\overline V$ and $v=\psi$ on that collar.
\end{lemma}
\begin{proof}
Let $|f_j|\leq M$. Take $\psi=A(|z|^2-R^2)+M+3$, with $A$ so large
that $\psi<-M-2$ on a neighborhood of $\overline V$. Since $\psi=M+3$
on $\partial U$, it exceeds $M+2$ on a sufficiently thin collar
$\{R-\delta_0<|z|<R\}$.

Use the regularized maximum $M_\epsilon$ constructed in
\cref{hessian:regularized-maximum}. Choose $0<\epsilon_j\downarrow0$
with $\epsilon_j<1$, and put $v_j=M_{\epsilon_j}(f_j,\psi)$.
That lemma shows that $v_j$ is smooth local $m$-subharmonic. The
monotonicity of $M_\epsilon$ in its arguments and in $\epsilon$ gives
$v_j\downarrow\max(f,\psi)=v$. Since $f_j-\psi>2>\epsilon_j$ near
$\overline V$ and $\psi-f_j>2>\epsilon_j$ on the collar, the two
asserted regions of exact equality follow. The bounds
\[
 \max(f_j,\psi)\leq v_j\leq\max(f_j,\psi)+\epsilon_j/2
\]
give a common bound. Since every $v_j$ equals $\psi$ on the collar,
all differences $v_k-v_j$ have support in the fixed compact set
$\{|z|\leq R-\delta_0\}$.
\end{proof}

\Needspace{14\baselineskip}
\begin{proposition}\label{hessian:smooth-capacity-convergence}\index{Quasicontinuity!Hessian capacity}
For the sequence $v_j\downarrow v$ of
\cref{hessian:boundary-localization} and $0\leq s\leq m$,
\[
 c_s\{v_k-v>\delta\}\longrightarrow0
                 \quad(k\to\infty)\qquad\text{for every }\delta>0.
\]
Consequently the pointwise limit of a uniformly bounded decreasing
sequence of smooth local $m$-subharmonic functions is quasicontinuous
on every compact subset of $U$, for each $c_s^*$ as well as $c_s$,
$0\leq s\leq m$.
\end{proposition}
\begin{proof}
The assertion for $s=0$ follows by bounded convergence for volume.
Induct on $s$, assuming it for $s-1$, and put $q=s-1$ and
$\theta=\omega^{n-s}$. For $j\geq k$ let $h=v_k-v_j\geq0$.
The support of $h$ lies in one fixed compact set $K_0\Subset U$.
The capacity bound \eqref{eq:hessian:smooth-capacity-bound} and the
induction hypothesis give
\begin{equation}\label{eq:hessian:lower-weight}
 \int h\,\ddc g_1\cdots\ddc g_q\omega^{n-q}
                          \longrightarrow0\quad(k\to\infty),
\end{equation}
uniformly in $j\geq k$ and in the smooth local $m$-subharmonic
$g_1,\ldots,g_q$ with a fixed common bound.
To see this, split at $h=\delta$. On $\{h\leq\delta\}$ the mass
bound of \cref{hessian:local-smooth-bounds} gives $C\delta$. On the
remaining set use $h\leq v_k-v$, the capacity bound with $s-1$ in
place of $s$, and the induction hypothesis. Let $k$ tend to
infinity before $\delta\downarrow0$. For $q=0$ this is the volume
estimate.

Consider the real smooth energy
\[
 E_s(w)=\int_U w(\ddc w)^s\theta,
 \qquad w_t=v_j+t h,\quad0\leq t\leq1.
\]
Each $w_t=(1-t)v_j+tv_k$ is smooth local $m$-subharmonic, with the
common bound of the $v_j$. The energies $E_s(v_j)$ are bounded: on
the fixed collar the integrand is fixed, and on the remaining compact
set \cref{hessian:local-smooth-bounds} gives a uniform mass bound.
Since $h$ vanishes near the boundary,
integration by parts and closedness of $\ddc w_t$ give the exact
identity
\begin{equation}\label{eq:hessian:energy-variation}
 \begin{split}
 \frac{\mathrm d}{\mathrm dt}E_s(w_t)
  &=(s+1)\int h(\ddc w_t)^s\theta\\
  &\quad+s\int h(\ddc w_t)^q
       \bigl(\mathrm dw_t\wedge\dc\theta
           -\dc w_t\wedge\mathrm d\theta+w_t\ddc\theta\bigr).
 \end{split}
\end{equation}
Here
\[
 \begin{aligned}
 \mathrm d\theta&=(n-s)\mathrm d\omega\wedge\omega^{n-s-1},\\
 \ddc\theta&=(n-s)\ddc\omega\wedge\omega^{n-s-1}\\
 &\quad+(n-s)(n-s-1)\mathrm d\omega\wedge\dc\omega
                                           \wedge\omega^{n-s-2},
 \end{aligned}
\]
and the analogous $\dc$ formula holds. A zero scalar coefficient
means that its term is absent. As $k\to\infty$, the two first-order
terms in \eqref{eq:hessian:energy-variation} tend to zero uniformly
in $t$ and $j\geq k$, by \cref{hessian:smooth-torsion}, the gradient
bound of \cref{hessian:local-smooth-bounds}, and
\eqref{eq:hessian:lower-weight}. The last term has the same
property by the absolute-form estimate at degree $q<m$ and
\eqref{eq:hessian:lower-weight}.
More explicitly, the first torsion integral is bounded by
\[
 C\left(\int h\,\mathrm dw_t\wedge\dc w_t
                     \wedge(\ddc w_t)^q\omega^{n-q-1}\right)^{1/2}
   \left(\int h(\ddc w_t)^q\omega^{n-q}\right)^{1/2}.
\]
The first factor is bounded since $h$ has a common bound and
\cref{hessian:local-smooth-bounds} applies; the second tends to
zero by \eqref{eq:hessian:lower-weight}. Conjugation handles the
other first-order term.
Integrating \eqref{eq:hessian:energy-variation} over $0\leq t\leq1$
therefore gives, for $j\geq k$,
\[
 E_s(v_k)-E_s(v_j)=(s+1)\Delta_{kj}+r_{kj},
 \quad
 \Delta_{kj}=\int_0^1\int h\Lambda_s(w_t)\,\mathrm dt\geq0,
 \quad |r_{kj}|\leq e_k,
\]
where $e_k\to0$ does not depend on $j$.
It follows that $E_s(v_j)\leq E_s(v_k)+e_k$ for $j\geq k$, so
$\limsup_jE_s(v_j)\leq E_s(v_k)+e_k$ for every $k$. Letting $k$ run
through a subsequence along which $E_s(v_k)$ tends to its lower
limit shows that the bounded sequence $E_s(v_j)$ converges. Hence
$\Delta_{kj}\to0$ as $k\to\infty$, uniformly for $j\geq k$.
Positivity of the binomial expansion
$\Lambda_s(w_t)\geq(1-t)^s \Lambda_s(v_j)$ then gives
\begin{equation}\label{eq:hessian:self-weight}
 \int h\Lambda_s(v_j)\longrightarrow0
             \quad(k\to\infty),\ \text{uniformly for }j\geq k.
\end{equation}

It remains to replace each Hessian factor by a capacity competitor.
For any smooth local $m$-subharmonic $a$ with $-1\leq a\leq0$, put
\[
 \begin{aligned}
 I_\ell&=\int h(\ddc v_j)^\ell(\ddc a)^{s-\ell}\theta,
   &&0\leq\ell\leq s,\\
 \eta_\ell&=(\ddc v_j)^\ell(\ddc a)^{s-\ell-1},
   &&0\leq\ell<s.
 \end{aligned}
\]
Stokes' theorem for the compactly supported form
$h\,\dc a\wedge\eta_\ell\theta$, with every metric derivative
kept, gives
\[
 I_\ell=-\int\mathrm dh\wedge\dc a\wedge\eta_\ell\theta
          +\int h\,\dc a\wedge\eta_\ell\mathrm d\theta.
\]
The second term is $o(1)$ as $k\to\infty$, uniformly in $a$ and
$j\geq k$. Indeed, apply the mixed form of
\cref{hessian:smooth-torsion} with $f=a$ and weight $h$; each subset
sum has a uniform gradient bound by \cref{hessian:local-smooth-bounds},
and its lower trace is controlled by \eqref{eq:hessian:lower-weight}.

The form $\eta_\ell\theta$ is positive of bidegree $(n-1,n-1)$:
a complementary nonnegative rank-one form completes a polarized
product of $s\leq m$ cone forms. Its associated Hermitian
gradient pairing is nonnegative, so Cauchy--Schwarz and the uniform
gradient bound for $a$ give
\[
 I_\ell\leq C F_\ell(h)^{1/2}+o(1),
 \qquad F_\ell(h)=\int\mathrm dh\wedge\dc h\wedge\eta_\ell\theta.
\]
Stokes' theorem for $h\,\dc h\wedge\eta_\ell\theta$ gives
\[
 F_\ell(h)= I_{\ell+1}
       -\int h\ddc v_k\wedge\eta_\ell\theta
       +\int h\,\dc h\wedge\eta_\ell\mathrm d\theta
                    \leq I_{\ell+1}+o(1).
\]
The dropped term is nonnegative. For the last term the same
polarized torsion estimate applies, with $f=h$ and weight $h$. Its
first gradient factor is uniformly bounded because $0\leq h\leq C$
and $h=v_k-v_j$: for every positive single-power gradient pairing,
the gradient energy of the difference is at most twice the sum of
those of $v_k$ and $v_j$. Its second factor tends to zero by
\eqref{eq:hessian:lower-weight}.
Starting with \eqref{eq:hessian:self-weight} at $I_s$, finite
backward induction gives $I_0\to0$ uniformly in $a$ and $j\geq k$.

Since $I_0=\int h\Lambda_s(a)$, Chebyshev's inequality now yields
$c_s\{v_k-v_j>\delta\}\leq\delta^{-1}\sup_a I_0$.
These sets increase to $\{v_k-v>\delta\}$ as $j\to\infty$.
Continuity of $c_s$ under increasing unions proves the induction.

For quasicontinuity, fix $0\leq s\leq m$ and $\epsilon>0$, and
choose $k_N$ so that the sets
$G_N=\{v_{k_N}-v>2^{-N}\}$ satisfy
$c_s(G_N)<\epsilon2^{-N}$, $N\geq1$.
They are open because $v$ is upper semicontinuous and $v_{k_N}$
is smooth. On $U\setminus\bigcup_NG_N$ these smooth functions
converge uniformly to the pointwise limit $v$. Hence $v$ is
continuous there. The union has capacity less than $\epsilon$.
Every compact subset of $U$ lies in a smaller concentric ball $V$,
and $v=f$ near $\overline V$ in \cref{hessian:boundary-localization};
this transfers the conclusion to $f$. Open sets have the same outer
capacity, which also proves the assertion for $c_s^*$.
\end{proof}

\Needspace{16\baselineskip}
\begin{lemma}\label{hessian:rough-coefficient-capacity}
Every bounded admissible potential on $X$ is locally quasicontinuous
for each smooth capacity $c_s^*$, $0\leq s\leq m$. For $0\leq p<m$,
each coefficient measure $\nu$ of the ordered $D_p$ satisfies
\[
 |\nu|(E)\leq C c_{m-1}^*(E),\qquad E\subset K\Subset U,
\]
where $C$ depends only on $K$, $U$, $\omega$ and a common bound for
the potentials. For $p=0$ the assertion concerns the volume measure,
the coefficient of $D_0=1$.
\end{lemma}
\begin{proof}
Let $u$ be bounded admissible. The global decreasing approximation
of \cref{hessian:bounded-decreasing-approximation} and
\cref{hessian:local-unshifted-correction} provide uniformly bounded
smooth local $m$-subharmonic functions decreasing to $u+\rho$ on $U$.
By \cref{hessian:smooth-capacity-convergence}, this pointwise limit,
formed with the canonical values of $u$, is quasicontinuous.
Subtracting the smooth $\rho$ preserves quasicontinuity.

For smooth admissible potentials $u_i$ write
$\ddc u_i=\ddc(u_i+\rho)-\ddc\rho$ and expand the product.
Each term has $q\leq p<m$ local cone factors and fixed smooth
complementary factors. The absolute-form estimate and
\eqref{eq:hessian:smooth-capacity-bound}, with $s=m-1$, bound
each coefficient variation by $C c_{m-1}$ on Borel sets.
The constants depend only on the common bound and local data.

Take the iterated limits of
\cref{hessian:bounded-recursive-currents} on a slightly larger
compact neighborhood of $K$. Lower semicontinuity of signed
measure variation on open sets preserves these estimates at each
limit. For any open $G\supset E$, intersect with a fixed open
neighborhood of $K$ having compact closure in $U$; this gives
$|\nu|(E)\leq Cc_{m-1}(G)$. Taking the infimum proves the outer
estimate.
\end{proof}

\begin{lemma}\label{hessian:weighted-measure-convergence}
Fix $0\leq s\leq m$ and the outer smooth capacity $c_s^*$ on $U$.
Let signed Radon measures $\nu_j\to\nu$ weakly have uniformly bounded
local variations, with all variations dominated on compact subsets
by $C c_s^*$. Suppose, for $1\leq i\leq\ell$, that bounded Borel
functions $b_i^j\downarrow b_i$ have one common bound, and that every
$b_i^j$ and every limit $b_i$ is locally quasicontinuous for this
capacity. For any continuous real $F$ on a compact box containing
their ranges,
\[
 F(b_1^j,\ldots,b_\ell^j)\nu_j
       \longrightarrow F(b_1,\ldots,b_\ell)\nu
\]
weakly as signed Radon measures.
\end{lemma}
\begin{proof}
Fix $\phi\in C_c(U)$, a compact $K\Subset U$ containing its
support, and $\epsilon>0$. For each of the countably many functions
$b_i^j$ and $b_i$, choose an exceptional open set for
quasicontinuity on $K$, with capacities summing to less than
$\epsilon$. Their union $G$ is open, $c_s(G)<\epsilon$, and every
function is continuous on the compact set $A=K\setminus G$.
Dini's theorem gives $b_i^j\to b_i$ uniformly on $A$ for each $i$.
Uniform continuity of $F$ on the compact box therefore gives
$f_j\to f$ uniformly on $A$, where
$f_j=F(b_1^j,\ldots,b_\ell^j)$ and $f=F(b_1,\ldots,b_\ell)$.
Put $M=\|F\|_\infty$. Extend $f|_A$ to a continuous $g$ on $U$
with $|g|\leq M$; if $A$ is empty, take $g=0$.
Then
\[
\begin{aligned}
 \left|\int\phi f_j\,\mathrm d\nu_j-
                 \int\phi f\,\mathrm d\nu\right|
 &\leq \|\phi\|_\infty\|f_j-f\|_A|\nu_j|(K)
       +\left|\int\phi g\,\mathrm d(\nu_j-\nu)\right|\\
 &\quad+2M\|\phi\|_\infty
            \bigl(|\nu_j|(K\cap G)+|\nu|(K\cap G)\bigr).
\end{aligned}
\]
For fixed $G$, the first two terms tend to zero as $j\to\infty$:
the first by uniform convergence on $A$, the second by weak
convergence applied to $\phi g\in C_c(U)$. Since
$|\nu_j|(K\cap G)\leq Cc_s(G)$, and similarly for $\nu$, the last
term is at most $4MC\|\phi\|_\infty\epsilon$. Let $\epsilon$ tend
to zero.
\end{proof}

\Needspace{20\baselineskip}
\begin{theorem}\label{hessian:joint-decreasing-products}\label{hessian:bounded-two-factor-convergence}\label{hessian:bounded-linear-weighted-convergence}\index{Hessian measure!simultaneous decreasing convergence}\index{Hessian measure!two bounded factors}\index{Hessian measure!weighted two-factor convergence}
For uniformly bounded admissible potentials $u_i^j\downarrow u_i$
on $X$, the currents $D_p(u_1^j,\ldots,u_p^j)$ converge to
$D_p(u_1,\ldots,u_p)$ in distributions for $p\leq m$, and weakly
with measure coefficients for $p<m$. These currents are symmetric
in their potentials. For $p<m$ the coefficient convergence also
holds after multiplication by a uniformly bounded decreasing
admissible weight. The positive shifted currents $Q_p$ converge
weakly for $p\leq m$; $Q_m$ converges as a positive Radon measure.
\end{theorem}
\begin{proof}
Work on a compact coordinate test set, and induct on $p$; the case
$p=0$ is trivial since $D_0=1$. Suppose the result holds for
$D_{p-1}$, which has measure coefficients since $p-1<m$.
\Cref{hessian:rough-coefficient-capacity,hessian:weighted-measure-convergence},
applied coefficientwise with $s=m-1$ and weight $u_p^j$, give weak
convergence of $u_p^jD_{p-1}(u_1^j,\ldots,u_{p-1}^j)$ to
$u_pD_{p-1}(u_1,\ldots,u_{p-1})$. Applying $\ddc$ proves
distributional convergence at degree $p$. For $p<m$, the uniform
coefficient bounds in \cref{hessian:bounded-recursive-currents}
upgrade this to weak convergence of measures, by smooth
approximation of compactly supported continuous tests.
\Cref{hessian:weighted-measure-convergence}, now with a decreasing
weight $a_j\downarrow a$, gives the additional weighted conclusion.

For each fixed list choose simultaneous smooth decreasing global
approximations. Their smooth raw Hessian wedges are symmetric.
The convergence just proved identifies their limit with $D_p$
in any ordering. This proves symmetry without exchanging the
iterated limits of \eqref{eq:hessian-bounded-iterated-limit}.

The subset expansion for $R_p$ uses fixed smooth powers of
$\omega$, so its distributional convergence and symmetry follow
from those of $D_I$. Its wedge $Q_p$ is positive with uniform
local masses by \cref{hessian:bounded-shifted-positive}.
Distributional convergence and these mass bounds imply weak
convergence as an order-zero positive current, including the
Radon measure conclusion at $p=m$. A partition of unity on a
finite coordinate cover proves the global assertions.
\end{proof}

The coefficient bound for the raw currents $D_p$, $p<m$, uses only
$c_{m-1}$. For the positive shifted products, $c_m$ controls the
traces in every degree $p\leq m$, including the top degree, where
$D_m$ itself has no coefficient bound. This allows a bounded
decreasing weight in every Hessian degree.

\Needspace{15\baselineskip}
\begin{lemma}\label{hessian:positive-capacity-domination}
On $K\Subset U$, the positive shifted currents
$Q_p(u_1,\ldots,u_p)$, $1\leq p\leq m$, for globally bounded
admissible potentials have trace and coefficient variations bounded
on Borel $E\subset K$ by $C c_m^*(E)$. The constants are uniform
when all potentials have a common bound.
\end{lemma}
\begin{proof}
First take smooth potentials, put $\sigma_i=\omega+\ddc u_i$, and
let $\rho$ be the fixed local correction of
\cref{hessian:local-unshifted-correction}, with
$\tau=\ddc\rho\geq2\omega$. Then
$\beta_i=\ddc(u_i+\rho)=\sigma_i+(\tau-\omega)$ lies in
$\overline{\Gamma_m(\omega)}$. Since $\tau-\omega\geq0$, mixed
positivity gives
\[
 0\leq\int_E Q_p\wedge\omega^{m-p}
    \leq\int_E\beta_1\cdots\beta_p\omega^{n-p}
    \leq C c_m(E),\qquad E\subset K.
\]
Here $Q_p\wedge\omega^{m-p}=\sigma_1\cdots\sigma_p\omega^{n-p}$.
The last bound is the domination of mixed traces by $c_m$ that
follows \eqref{eq:hessian:smooth-capacity-bound}, applied with $s=m$
to the corrected potentials $u_i+\rho$, which are bounded on $U$.
Positivity of $Q_p$ bounds each of its
coefficient variations on the compact coordinate set by a fixed
multiple of its trace measure.

For bounded potentials choose uniformly bounded smooth decreasing
approximations by \cref{hessian:bounded-decreasing-approximation}. The joint
convergence in \cref{hessian:joint-decreasing-products} identifies
their positive-current limit with the previously constructed $Q_p$.
On a compact neighborhood of $K$ inside $U$, lower semicontinuity
on open sets preserves trace and coefficient variation estimates.
For an arbitrary open $G\supset E$ intersect with a fixed relatively
compact neighborhood of $K$ before taking this limit; its capacity
is at most $c_m(G)$. Therefore every such trace or variation measure
$\mu$ satisfies $\mu(E)\leq Cc_m(G)$. Infimize over $G$ to obtain
the asserted outer-open bound. The constants of the smooth estimates
are independent of the approximation index.
\end{proof}

\Needspace{15\baselineskip}
\begin{theorem}\label{hessian:weighted-positive-products}\index{Hessian measure!weighted positive products}
Suppose $a_j\downarrow a$ and $u_i^j\downarrow u_i$ are bounded
admissible potentials on $X$ with common bounds. Then,
for $1\leq p\leq m$,
\[
 a_j Q_p(u_1^j,\ldots,u_p^j)
        \longrightarrow a Q_p(u_1,\ldots,u_p)
\]
weakly as currents with measure coefficients. At $p=m$ this is weak
convergence of weighted Radon measures.
\end{theorem}
\begin{proof}
Work on a compact coordinate test set. By
\cref{hessian:joint-decreasing-products}, every coefficient measure
of the positive shifted currents converges weakly and has a uniform
local variation bound. \Cref{hessian:positive-capacity-domination}
strengthens this to domination by $C c_m^*$ on Borel sets.
\Cref{hessian:rough-coefficient-capacity} makes the weights
$a_j$ and $a$ quasicontinuous for this capacity. Apply
\cref{hessian:weighted-measure-convergence} with $s=m$ to each
coefficient, and then use a finite partition of unity. At $p=m$,
$Q_m$ is a positive Radon measure, so this also gives the stated
weighted measure convergence.
\end{proof}

\section{Locality and maxima of bounded potentials}
\label{sec:hessian-bounded-locality}

Where one bounded admissible potential is strictly larger than another,
the maximum has the same positive Hessian current as the larger
potential. These strict comparison sets need not be open.
Quasicontinuity gives cutoffs whose products with the approximating
currents converge, allowing smooth locality to pass to these sets.

\Needspace{15\baselineskip}
\begin{proposition}\label{hessian:bounded-strict-locality}\index{Hessian measure!locality}
Let $u,v$ be bounded admissible, put $w=\max\{u,v\}$, and fix
$1\leq p\leq m$. For any bounded admissible $t_2,\ldots,t_p$,
\[
 \mathbf1_{\{u>v\}}Q_p(w,t_2,\ldots,t_p)
       =\mathbf1_{\{u>v\}}Q_p(u,t_2,\ldots,t_p).
\]
The equality is coefficientwise as currents with measure
coefficients. The corresponding equality on $\{v>u\}$ follows
by exchanging $u,v$.
\end{proposition}
\begin{proof}
Choose smooth uniformly bounded decreasing approximations of all
the potentials by \cref{hessian:bounded-decreasing-approximation}.
Construct $w_j$ by \cref{hessian:regularized-maximum}. Work on
one compact coordinate test set. Fix $\delta>0$ and a smooth
$\chi_\delta:\mathbb R\to[0,1]$ that vanishes for
$s\leq\delta/2$ and equals one for $s\geq\delta$.
For $j$ large enough that $\epsilon_j<\delta/2$, ordinary smooth
locality gives
\[
 \chi_\delta(u_j-v_j)Q_p(w_j,t_2^j,\ldots,t_p^j)
       =\chi_\delta(u_j-v_j)Q_p(u_j,t_2^j,\ldots,t_p^j).
\]
Indeed, wherever the weight is nonzero, $w_j=u_j$ on an open
neighborhood. The equality elsewhere follows from its zero weight.

Both sequences of positive currents converge by
\cref{hessian:joint-decreasing-products}; every coefficient variation
is uniformly dominated by $C c_m^*$ by
\cref{hessian:positive-capacity-domination}. The potentials and their
limits are quasicontinuous by \cref{hessian:rough-coefficient-capacity}.
Apply \cref{hessian:weighted-measure-convergence} with $s=m$ and
$F(x,y)=\chi_\delta(x-y)$ to the real and imaginary parts of
each coefficient. Passing to the limit gives the same identity
with $u,v,w,t_2,\ldots,t_p$.
Choose the cutoffs so that $\chi_\delta(s)\to\mathbf1_{\{s>0\}}$
pointwise as $\delta\downarrow0$. Dominated convergence against
the fixed finite coefficient variations on the test set gives
the asserted restriction. A partition of unity proves it globally.
\end{proof}

On the contact set, a positive constant shift gives a strict comparison
before the limit is taken. Decreasing the shift then recovers that
set by dominated convergence for a fixed measure.

\Needspace{16\baselineskip}
\begin{theorem}\label{hessian:bounded-maximum-measure}\index{Hessian measure!maximum inequality}
Let $u,v$ be bounded admissible. Fix $1\leq r\leq m$ and bounded admissible
$z_1,\ldots,z_{m-r}$. Define the positive Radon measure
\[
 \mu_r(t)=Q_m(\underbrace{t,\ldots,t}_{r},z_1,\ldots,z_{m-r}).
\]
For $w=\max\{u,v\}$,
\[
 \mu_r(w)\geq\mathbf1_{\{u\geq v\}}\mu_r(u)
                  +\mathbf1_{\{u<v\}}\mu_r(v).
\]
The inequality with $u,v$ exchanged also holds.
\end{theorem}
\begin{proof}
Symmetry and \cref{hessian:bounded-strict-locality}, applied to the $r$ slots, give
$\mu_r(w)=\mu_r(v)$ on $\{u<v\}$ and $\mu_r(w)=\mu_r(u)$ on $\{u>v\}$.
To include the contact set, put $w_\epsilon=\max\{u+\epsilon,v\}$.
Strict locality and invariance under addition of constants give
\[
 \mu_r(w_\epsilon)\geq
                   \mathbf1_{\{u+\epsilon>v\}}\mu_r(u).
\]
The invariance follows from the smooth products and their joint
decreasing limits. As $\epsilon\downarrow0$,
$w_\epsilon\downarrow w$, so the left side converges weakly to
$\mu_r(w)$. Against any nonnegative continuous compactly supported
test function, the right side converges by dominated convergence
for the fixed measure $\mu_r(u)$ to
$\mathbf1_{\{u\geq v\}}\mu_r(u)$. This proves the measure inequality
on the contact set as well as on $\{u>v\}$. On its disjoint
complement $\{u<v\}$, the already proved equality supplies
$\mu_r(v)$. Restricting the positive measures to these two Borel sets
and adding gives the assertion. Exchanging $u,v$ gives its other
form.
\end{proof}

\section{Hessian mass and global comparison}
\label{sec:hessian-mass-comparison}

The maximum inequality compares measures on complementary Borel sets.
Integrating it leaves a difference of total masses. The recursive
currents express this difference through the complex Hessians of the
background powers, extending the mass calculation of
\cref{chap:mass} to bounded admissible potentials.

For bounded admissible $t_1,\ldots,t_m$, put
\[
 \mathcal M_m(t_1,\ldots,t_m)=\int_XQ_m(t_1,\ldots,t_m),
 \qquad H_m(t)=Q_m(t,\ldots,t).
\]
Here $H_m(t)$ is a finite positive Radon measure and
$\mathcal M_m$ is the total mass of a mixed positive measure, by
\cref{hessian:bounded-shifted-positive}. For a nonempty subset
$I=\{i_1<\cdots<i_q\}\subset\{1,\ldots,m\}$, put
$i(I)=i_q$ and $I^-=I\setminus\{i_q\}$. Use $D_{I^-}$ for the
recursive current with those ordered inputs, with $D_\varnothing=1$.

\Needspace{17\baselineskip}
\begin{lemma}\label{hessian:bounded-mass-expansion}\index{Hessian measure!total mass}
For every bounded admissible list,
\begin{equation}\label{eq:hessian-bounded-mass-expansion}
 \mathcal M_m(t_1,\ldots,t_m)
 =\int_X\omega^n+
 \sum_{\varnothing\ne I\subset\{1,\ldots,m\}}
 \int_X t_{i(I)}D_{I^-}\wedge\ddc(\omega^{n-|I|}).
\end{equation}
Every integral on the right is an integral of bounded canonical
values against signed measure coefficients.
\end{lemma}
\begin{proof}
Expand $Q_m$ using \eqref{eq:hessian-bounded-shifted-current}.
The term for $I$ pairs $D_I$ with the smooth form
$\omega^{n-|I|}$. By \eqref{eq:hessian-bounded-recursion},
\[
 \langle D_I,\omega^{n-|I|}\rangle
 =\langle t_{i(I)}D_{I^-},\ddc(\omega^{n-|I|})\rangle.
\]
Since $|I^-|\leq m-1$, \cref{hessian:bounded-recursive-currents}
gives measure coefficients for $D_{I^-}$; their bounded product
with $t_{i(I)}$ is well defined. Thus the displayed pairing is
exactly the integral in the sum, and is finite on compact $X$.
The empty subset contributes $\int_X\omega^n$.
\end{proof}

Let $\mu_r$ be the repeated-slot mixed measure from
\cref{hessian:bounded-maximum-measure}, for fixed $r$ and
$z_1,\ldots,z_{m-r}$.
The total-mass difference below therefore has the explicit expression
\eqref{eq:hessian-bounded-mass-expansion} for the two indicated lists.

\Needspace{11\baselineskip}
\begin{proposition}\label{hessian:comparison-mass-defect}\index{Comparison principle!Hessian mass defect}
If $u,v$ are bounded admissible and $w=\max\{u,v\}$, then
\[
 \int_{\{u<v\}}\mathrm d\mu_r(v)
 \leq\int_{\{u<v\}}\mathrm d\mu_r(u)
                +\int_X\mathrm d\mu_r(w)-\int_X\mathrm d\mu_r(u).
\]
\end{proposition}
\begin{proof}
Set $E=\{u<v\}$. The contact-set version of
\cref{hessian:bounded-maximum-measure} gives
\[
 \mu_r(w)\geq\mathbf1_{X\setminus E}\mu_r(u)
                     +\mathbf1_E\mu_r(v).
\]
All these measures are finite and positive. Integrating on $X$ and
using $\mu_r(u)(X\setminus E)=\mu_r(u)(X)-\mu_r(u)(E)$ gives
exactly the assertion. No boundary or closed-background hypothesis
enters this calculation.
\end{proof}

The unmodified global comparison principle holds precisely when the
powers occurring in this mass formula are $\ddc$-closed. Only the
powers $n-m,\ldots,n-1$ occur at Hessian degree $m$.

\Needspace{24\baselineskip}
\begin{theorem}\label{hessian:global-comparison-characterization}\index{Hessian measure!invariant mass}\index{Comparison principle!Hessian power identities}
The following conditions are equivalent:
\begin{enumerate}
\item $\ddc(\omega^{n-q})=0$ for every $1\leq q\leq m$;
\item $\int_XH_m(u)$ is independent of the smooth admissible
      potential $u$;
\item $\mathcal M_m(t_1,\ldots,t_m)=\int_X\omega^n$ for every
      bounded admissible list;
\item every pair of bounded admissible $u,v$ satisfies
\[
 \int_{\{u<v\}}H_m(v)\leq\int_{\{u<v\}}H_m(u).
\]
\end{enumerate}
Under these conditions, for every $1\leq r\leq m$ and every fixed bounded
admissible list $z_1,\ldots,z_{m-r}$, one also has
\[
 \int_{\{u<v\}}\mathrm d\mu_r(v)
           \leq\int_{\{u<v\}}\mathrm d\mu_r(u).
\]
\end{theorem}
\begin{proof}
The first condition implies the third by
\eqref{eq:hessian-bounded-mass-expansion}. The third implies the
second by taking all inputs equal to a smooth potential.

To prove that the second implies the first, fix any smooth real
$f$. For all sufficiently small real $t$ of either sign,
$\omega+t\ddc f$ is positive, so $tf$ is strictly admissible.
The polynomial
\[
 \int_X(\omega+t\ddc f)^m\wedge\omega^{n-m}
 =\sum_{q=0}^m\binom mq t^q
              \int_X(\ddc f)^q\wedge\omega^{n-q}
\]
is constant on an interval. Consequently every coefficient with
$q\geq1$ vanishes. Polarizing that coefficient by replacing $f$
with $\sum_{a=1}^q s_af_a$ gives
\[
 \int_X\ddc f_1\wedge\cdots\wedge\ddc f_q
                        \wedge\omega^{n-q}=0
\]
for all smooth real $f_a$. Integration by parts in the first
factor, using \cref{lem:ibp}, gives
\[
 \int_Xf_1\ddc f_2\wedge\cdots\wedge\ddc f_q
                         \wedge\ddc(\omega^{n-q})=0.
\]
Since $f_1$ is arbitrary, the smooth top-degree form paired with
it vanishes pointwise. At each point the other complex Hessians
can be prescribed by real quadratic functions times a cutoff.
Complexifying this multilinear identity gives all coordinate
$(q-1,q-1)$-forms, as in the proof of
\cref{thm:mass-invariance-characterization}. They detect every
coefficient of $\ddc(\omega^{n-q})$, so that form vanishes. When
$q=1$, there are no other Hessians and the same conclusion follows
directly from the arbitrariness of $f_1$.

Under the third condition, $\mu_r(w)$ and $\mu_r(u)$ have the same
total mass. \Cref{hessian:comparison-mass-defect} proves the mixed
comparison assertion, and $r=m$ gives the fourth condition.
Finally, suppose the fourth condition holds and fix smooth admissible
$u,v$. Choose $C$ so large that $u<v+C$ everywhere.
Constant invariance gives $H_m(v+C)=H_m(v)$, so comparison gives
$\int_XH_m(v)\leq\int_XH_m(u)$. Exchanging $u,v$ gives equality,
which is the second condition.
\end{proof}

The same varying background that obstructs Monge--Amp\`ere comparison
also obstructs comparison at each intermediate Hessian degree.

\Needspace{20\baselineskip}
\begin{example}\label{hessian:torus-comparison-obstruction}\index{Torus!Hessian comparison obstruction}
Let $X=(\mathbb C/(\mathbb Z+\mathrm i\mathbb Z))^n$, with
$z_j=x_j+\mathrm i y_j$, and put
\[
 \beta_j=\frac{\mathrm i}{2\pi}\,\mathrm dz_j\wedge\mathrm d\bar z_j,
 \qquad b=2+\cos(2\pi x_1),
 \qquad\omega=\beta_1+b\beta_2+\sum_{j=3}^n\beta_j.
\]
For $|t|<\pi^{-2}$ the potential $u_t=t\cos(2\pi x_1)$ is smooth
strictly admissible, and, with $B_n=\int_X\beta_1\wedge\cdots\wedge\beta_n$,
\[
 \int_X H_m(u_t)=n!\left(2-\frac{m\pi^2t}{2n}\right)B_n.
\]
Thus every $-\pi^{-2}<t<0$, together with a constant $C>|t|$,
gives a pair $u=0$, $v=u_t+C$ that violates the unmodified global
comparison inequality.
\end{example}
\begin{proof}
As in \cref{ex:torus-mass},
$\ddc u_t=-\pi^2t\cos(2\pi x_1)\beta_1$.
The coefficient of $\beta_1$ in $\omega+\ddc u_t$ is strictly
positive in the stated range, and all the other coefficients are
positive. Since $(\ddc u_t)^2=0$, expansion gives
\[
 H_m(u_t)=\omega^n+m\ddc u_t\wedge\omega^{n-1}
 =n!b\left(1-\frac mn\pi^2t\cos(2\pi x_1)\right)
                         \beta_1\wedge\cdots\wedge\beta_n.
\]
The period averages of $b$ and $b\cos(2\pi x_1)$ are respectively
$2$ and $1/2$, proving the mass formula. For $t<0$ this mass is
strictly larger than $\int_XH_m(0)=2n!B_n$. Meanwhile $C>|t|$
gives $\{0<u_t+C\}=X$, so the comparison inequality for this pair
would assert the opposite ordering of these masses.
\end{proof}

\subsection{Domination and proper equations}
\label{subsec:hessian-bounded-domination}

The positive background term in a convex rescaling makes a strict
comparison set visible to volume. Under the power identities of
\cref{hessian:global-comparison-characterization}, the global comparison
principle then yields domination for bounded potentials. Canonical
representatives turn the resulting almost everywhere order into
pointwise order.

\Needspace{14\baselineskip}
\begin{theorem}\label{hessian:strict-bounded-domination}\index{Domination principle!bounded Hessian potentials}
Assume $\ddc(\omega^{n-q})=0$ for $1\leq q\leq m$. Let $u,v$
be bounded admissible. If, for some $0\leq c<1$,
\[
 \mathbf1_{\{u<v\}}H_m(u)
                   \leq c\mathbf1_{\{u<v\}}H_m(v),
\]
then $u\geq v$ everywhere. In particular, this conclusion holds
if $H_m(u)(\{u<v\})=0$.
\end{theorem}
\begin{proof}
Choose $B>\|v\|_\infty$ and, for $0<\alpha<1$ with $\alpha^m>c$,
put $v_\alpha=\alpha v-(1-\alpha)B$. Then $v_\alpha\leq v$ and
$E_\alpha=\{u<v_\alpha\}\subset\{u<v\}$.
By \cref{hessian:global-comparison-characterization} and the
assumed restriction inequality,
\[
 \int_{E_\alpha}H_m(v_\alpha)
        \leq\int_{E_\alpha}H_m(u)
        \leq c\int_{E_\alpha}H_m(v).
\]
On the other hand, \cref{hessian:convex-background-term} gives
\[
 (\alpha^m-c)\int_{E_\alpha}H_m(v)
             +(1-\alpha)^m\int_{E_\alpha}\omega^n\leq0.
\]
Both summands are nonnegative, and the coefficient of the volume
is positive. Thus $E_\alpha$ has zero volume. Applying
\cref{hessian:canonical-order} gives $u\geq v_\alpha$ pointwise.
Let $\alpha$ increase to one. The boundedness of $v$ gives
$v_\alpha\to v$ uniformly, so $u\geq v$ everywhere. The final
assertion is the case $c=0$.
\end{proof}

\Needspace{13\baselineskip}
\begin{corollary}\label{hessian:proper-bounded-uniqueness}\index{Hessian equation!bounded proper uniqueness}
Assume $\ddc(\omega^{n-q})=0$ for $1\leq q\leq m$ and fix $a>0$.
If bounded admissible $u,v$ satisfy
\[
 e^{-av}H_m(v)\geq e^{-au}H_m(u),
\]
then $u\geq v$. Consequently, for any fixed finite positive Radon
measure $\nu$, the equation $H_m(u)=e^{au}\nu$ has at most one
bounded admissible solution.
\end{corollary}
\begin{proof}
For $\delta>0$, put $E_\delta=\{u<v-\delta\}$. Multiplying the
assumed measure inequality by the positive bounded Borel function
$e^{au}$ gives on $E_\delta$
\[
 H_m(u)\leq e^{a(u-v)}H_m(v)
                 \leq e^{-a\delta}H_m(v-\delta).
\]
The last identity uses invariance under constants. Since
$e^{-a\delta}<1$, \cref{hessian:strict-bounded-domination} applied
to $u,v-\delta$ gives $u\geq v-\delta$ everywhere. Letting
$\delta\downarrow0$ gives $u\geq v$. Two solutions of the stated
equation have equal weighted measures $e^{-au}H_m(u)=\nu$.
Applying the order assertion in both directions proves uniqueness.
\end{proof}

A strictly increasing exponent also controls changes in logarithmic
data for a fixed Radon right-hand measure.

\Needspace{20\baselineskip}
\begin{proposition}\label{hessian:proper-bounded-stability}\index{Hessian equation!bounded data stability}
Assume $\ddc(\omega^{n-q})=0$ for $1\leq q\leq m$. Let $A<B$ and $f$ be
smooth near $X\times[A,B]$, with $\partial_t f\geq\kappa>0$
there, and let $\nu$ be a finite positive Radon measure.
Suppose bounded admissible $u,v$ have values in $[A,B]$ and satisfy
\[
 H_m(u)=e^{f(x,u)+h_1}\nu,\qquad
 H_m(v)=e^{f(x,v)+h_2}\nu,
\]
where $h_1,h_2$ are bounded Borel functions. Then
\[
 \|u-v\|_\infty
            \leq\kappa^{-1}\|h_1-h_2\|_{L^\infty(\nu)}.
\]
\end{proposition}
\begin{proof}
Put $K=\|h_1-h_2\|_{L^\infty(\nu)}$ and fix $s>K/\kappa$.
On $E_s=\{v+s<u\}$, the actual values $u,v$ lie in $[A,B]$,
so $f(x,u)-f(x,v)\geq\kappa s$. The equations and the essential
bound on $h_2-h_1$ therefore give, as restricted measures,
\[
 \mathbf1_{E_s}H_m(v+s)
 =\mathbf1_{E_s}H_m(v)
 \leq e^{K-\kappa s}\mathbf1_{E_s}H_m(u).
\]
The constant is smaller than one. Applying
\cref{hessian:strict-bounded-domination} to $v+s,u$ gives
$u\leq v+s$ pointwise. Let $s\downarrow K/\kappa$ and exchange
$u,v$ to obtain the norm bound. The exponent is evaluated only
at $u,v$; the shifted comparison potential need not take its
values in $[A,B]$.
\end{proof}

\Needspace{14\baselineskip}
\begin{corollary}\label{hessian:proper-bounded-smoothness}\index{Hessian equation!bounded proper smoothness}
Assume $\ddc(\omega^{n-q})=0$ for $1\leq q\leq m$.
Under the hypotheses of \cref{hessian:proper-variable-existence},
every bounded admissible solution with values in $[A,B]$ is its
smooth solution. In particular, for every $a>0$ and smooth real
$H$, the equation $H_m(u)=e^{au+H}\omega^n$ has exactly one bounded
admissible solution, and that solution is smooth and strictly
admissible.
\end{corollary}
\begin{proof}
For a continuous admissible potential, smooth decreasing admissible
approximations converge uniformly on compact $X$ by Dini's theorem.
Thus \cref{hessian:joint-decreasing-products,hessian:continuous-products}
define the same Hessian measure on the continuous class.
The smooth solution of \cref{hessian:proper-variable-existence}
therefore solves the bounded equation as well. Apply
\cref{hessian:proper-bounded-stability} with $\nu=\omega^n$ and
$h_1=h_2=0$ to identify any other bounded solution in $[A,B]$
with it. For the final assertion, smooth existence is already
\cref{hessian:smooth-proper-existence}, and
\cref{hessian:proper-bounded-uniqueness} with
$\nu=e^H\omega^n$ gives uniqueness among all bounded solutions.
\end{proof}

